% Mechanically extracted after \mainmatter from the frozen standalone master.

\markboth{Exposé XIII. Regular Elements}{Exposé XIII. Regular Elements}
% BEGIN inputs/XIII/source/tex/components/01_expose_XIII_opening.tex
% Authority: Expo13.pdf, local pp. 1--3 / combined pp. 805--807.
% Stops before: “Nous supposons par la suite que S est le spectre...”

\chapter*{Exposé XIII\\Regular Elements of Algebraic Groups and Lie Algebras}
\addcontentsline{toc}{chapter}{Exposé XIII. Regular Elements of Algebraic Groups and Lie Algebras}
\SGARefTarget{sga3:xiii:expose}

\begin{center}
\textit{by A. Grothendieck}
\end{center}

\section{An auxiliary lemma on varieties with operators}
\SGARefTarget{sga3:xiii:section:1}

Let \(S\) be a prescheme, let \(G\) be an \(S\)-group prescheme acting
on the left on an \(S\)-prescheme \(V\), let \(W\) be a closed
\(S\)-subprescheme of \(V\), and let \(N\) be its stabilizer in \(G\).
Thus \(N\) is the subgroup of \(G\) whose points with values in an
\(S'\) over \(S\) are the elements \(g\in G(S')\) such that
\[
  g\cdot W_{S'}=W_{S'}.
\]
We endow \((\mathrm{Sch}/S)\) with the faithfully flat
quasi-compact topology and identify \(G\), \(V\), and \(W\) with the
corresponding sheaves.\footnote{The French source prints \(M\) in this
list; the surrounding sentence and all subsequent uses require \(W\).}
We shall therefore reason in the category of sheaves on
\((\mathrm{Sch}/S)\); throughout \SGARefLink{sga3:xiii:section:1}{this section}, ``locally'' refers to
the topology just specified on \((\mathrm{Sch}/S)\).

Since \(N\) is a sheaf, consider the quotient sheaf \(G/N\). It is
immediately seen to be isomorphic to the following functor: to every
\(S'\) over \(S\), it assigns the set of subsheaves \(W'\) of
\(V_{S'}\) that are locally conjugate to \(W_{S'}\) under \(G\).
Let \(X\) be the subsheaf of \(G/N\times_S V\) whose value on an
\(S'\) over \(S\) is the set of pairs \((W',v)\), where \(W'\) is as
above and \(v\) is a section of \(W'\) over \(S'\) (hence also a
section of \(V_{S'}\) over \(S'\)). Let \(Z\) be the inverse image of
\(X\) in \(G\times_S V\). We then have \SGARefLink{sga3:xiii:equation:x}{the cartesian diagram}
\SGARefTarget{sga3:xiii:diagram:x}
\[
\tag{x}\label{sga3:xiii:equation:x}
\begin{tikzcd}[column sep=large,row sep=large]
Z \arrow[r] \arrow[d,hook] & X \arrow[d,hook,"i"] \\
G\times_S V \arrow[r] & G/N\times_S V ,
\end{tikzcd}
\]
where \(i\) is the canonical immersion, while the lower horizontal
arrow is induced by the canonical morphism \(G\to G/N\).

This morphism sends \(g\in G(S')\) to the subsheaf
\(g\cdot W_{S'}\) of \(V_{S'}\). Consequently, \(Z(S')\) is the set
of pairs \((g,v)\in G(S')\times V(S')\) such that
\(v\in (g\cdot W_{S'})(S')\). It follows that \(Z\) is isomorphic to
the sheaf \(G\times_S W\), via the isomorphism
\[
  G\times_S W \xrightarrow{\sim} Z,\qquad
  (g,w)\longmapsto(g,g\cdot w).
\]
Thus \SGARefLink{sga3:xiii:equation:x}{the preceding cartesian diagram} yields \SGARefLink{sga3:xiii:equation:xx}{the cartesian diagram}
\SGARefTarget{sga3:xiii:diagram:xx}
\[
\tag{xx}\label{sga3:xiii:equation:xx}
\begin{tikzcd}[column sep=large,row sep=large]
G\times_S W \arrow[r,"q"] \arrow[d,"\lambda"'] &
  X \arrow[d,"i"] \\
G\times_S V \arrow[r] & G/N\times_S V
\end{tikzcd}
\]
where
\[
  \lambda(g,w)=(g,g\cdot w),\qquad
  q(g,w)=(\bar g,g\cdot w),
\]
where \(\bar g\) denotes the image of \(g\) under the canonical map
\(G(S')\to(G/N)(S')\).

Diagram \(\SGARefLink{sga3:xiii:equation:x}{\mathrm{(x)}}\) also shows that \(Z\to X\) makes \(Z\) a
principal bundle over \(X\), with \(N\) acting on the right by
\[
  (g,v)\cdot n=(gn,v).
\]
Accordingly, in diagram \(\SGARefLink{sga3:xiii:equation:xx}{\mathrm{(xx)}}\), the morphism
\(q:G\times_S W\to X\) makes \(G\times_S W\) a principal bundle over
\(X\), with right \(N\)-action
\[
  (g,w)\cdot n=(gn,n^{-1}\cdot w).
\]

The principal morphisms just introduced are summarized by
\SGARefTarget{sga3:xiii:diagram:d}
\[
\tag{D}\label{sga3:xiii:equation:d}
\begin{tikzcd}[column sep=large,row sep=large]
& G\times_S W \arrow[d,"q"] \arrow[dr,"\varphi"] & \\
G/N &
X \arrow[l,"p"] \arrow[r,"\psi"] \arrow[d,hook,"i"] &
V \\
& G/N\times_S V
  \arrow[ul,dashed,"\operatorname{pr}_1"]
  \arrow[ur,dashed,"\operatorname{pr}_2"'] &
\end{tikzcd}
\]
where
\[
  \psi=\operatorname{pr}_2\circ i,\qquad
  \varphi=\psi\circ q,\qquad
  \varphi(g,w)=g\cdot w.
\]

If \(v\) is a section of \(V\) over \(S\), let \(X_v\) be its inverse
image in \(X\) under \(\psi\). For every \(S'\) over \(S\),
\(X_v(S')\) is the set of subsheaves \(W'\) of \(V_{S'}\) that are
locally conjugate to \(W_{S'}\) under \(G\) and contain the section
\(v_{S'}\) of \(V_{S'}\). The inverse image of \(v\) in
\(G\times_S W\) under \(\varphi\) is isomorphic to the subsheaf
\(M_v\) of \(G\) defined by
\[
  M_v(S')=
  \bigl\{g\in G(S')\mid v_{S'}\in(g\cdot W_{S'})(S')\bigr\}
  =
  \bigl\{g\in G(S')\mid g^{-1}v_{S'}\in W(S')\bigr\}.
\]
If \(v\) is a section of \(W\), and not merely of \(V\), then
\(M_v\) evidently contains \(N\).

None of these explicit descriptions used the representability of
\(G\), \(V\), or \(W\) (nor even the fact that the site in question is
defined in terms of preschemes). Suppose now, however, that \(N\) is
representable, faithfully flat, and quasi-compact over \(S\), and that
\(G/N\) is representable. When \(S\) is the spectrum of a field and
\(G\) is of finite type over that field, this hypothesis is necessarily
satisfied (VIII~6 and VI\textsubscript{B}.11.18). Using the cartesian
diagram \(\SGARefLink{sga3:xiii:equation:xx}{\mathrm{(xx)}}\), faithfully flat quasi-compact descent, and
the fact that \(Z\to G\times_S V\) is a closed immersion, one sees that
\(X\) is representable: it is obtained by descending the closed
subprescheme \(Z\) of \(G\times_S V\) along the faithfully flat
quasi-compact morphism
\[
  G\times_S V\longrightarrow G/N\times_S V.
\]
Thus diagram \(\SGARefLink{sga3:xiii:equation:d}{\mathrm{(D)}}\) is a diagram of morphisms of
preschemes over \(S\).

% END inputs/XIII/source/tex/components/01_expose_XIII_opening.tex

% BEGIN inputs/XIII/source/tex/components/02_expose_XIII_section1_remainder_and_section2.tex
% Authority: Expo13.pdf, from “Nous supposons par la suite que S est le
% spectre d’un corps k...” through the end of §2, immediately before §3.
% PDF local pp. 3--14 (requested boundary: local p. 3 / combined p. 807
% through local p. 14); printed pp. 251--269.

We henceforth assume that \(S\) is the spectrum of a field \(k\), and
that \(G\), \(V\), and \(W\) are of finite type over \(k\). Let
\(\mathfrak n\) be the Lie algebra of \(N\). Then
\(\dim N\leq\operatorname{rank}_k\mathfrak n\), with equality if and
only if \(N\) is smooth over \(k\) (Exposé VI\footnote{Editor's note in
the French source: since this reference could not be identified, see
Theorem II.5.2.1 of M.~Demazure and P.~Gabriel,
\emph{Groupes algébriques I}, Masson (1970).}). Let \(a\in W(k)\), and
consider the subscheme \(M_a\) of \(G\) defined above; it contains
\(N\) and is isomorphic to \(\varphi^{-1}(a)\). We denote by
\(\mathfrak m_a\) its Zariski tangent space at the identity element
\(e\) of \(N\). Thus\footnote{The French source omits
\(\operatorname{rank}_k\) before \(\mathfrak m_a\) in \SGARefLink{sga3:xiii:equation:1}{(1)}; it is
restored here, as required by the surrounding text and \SGARefLink{sga3:xiii:lemma:1:1}{Lemma 1.1}.}
\SGARefTarget{sga3:xiii:equation:1}
\[
\tag{1}
  \mathfrak n\subset\mathfrak m_a,\qquad
  \dim N\leq\operatorname{rank}_k\mathfrak n
  \leq\operatorname{rank}_k\mathfrak m_a.
\]

\medskip
\noindent\SGARefTarget{sga3:xiii:lemma:1:1}\textbf{Lemma 1.1.} ---
\textit{With the preceding notation:}

\begin{enumerate}
\item[\textit{a)}] \SGARefTarget{sga3:xiii:lemma:1:1:item:a} Consider the following conditions:
  \begin{enumerate}
  \item[(i)] \SGARefTarget{sga3:xiii:lemma:1:1:item:i}
    \(\mathfrak n=\mathfrak m_a\), and \(N\) is smooth over \(k\).
  \item[(i bis)] \SGARefTarget{sga3:xiii:lemma:1:1:item:i-bis}
    \(\dim N=\operatorname{rank}_k\mathfrak m_a\).
  \item[(ii)] \SGARefTarget{sga3:xiii:lemma:1:1:item:ii}
    The morphism \(\psi:X\to V\) is unramified at
    \((\bar e,a)\).
  \item[(iii)] \SGARefTarget{sga3:xiii:lemma:1:1:item:iii}
    \(M_a\) and \(N\) coincide in a neighborhood of \(e\).
  \end{enumerate}
  Then
  \[
    \SGARefLink{sga3:xiii:lemma:1:1:item:i}{\mathrm{(i)}}\Longleftrightarrow\SGARefLink{sga3:xiii:lemma:1:1:item:i-bis}{\mathrm{(i\ bis)}}
    \Longrightarrow\SGARefLink{sga3:xiii:lemma:1:1:item:ii}{\mathrm{(ii)}}
    \Longleftrightarrow\SGARefLink{sga3:xiii:lemma:1:1:item:iii}{\mathrm{(iii)}}.
  \]
\item[\textit{b)}] \SGARefTarget{sga3:xiii:lemma:1:1:item:b}
  Suppose that \(\varphi:G\times_S W\to V\) is smooth at \((e,a)\).
  Then \(M_a\) is smooth over \(k\) at \(e\), and \(\psi\) is smooth
  at \((\bar e,a)\).
\end{enumerate}

\noindent\textit{Proof.}
\SGARefTarget{sga3:xiii:lemma:1:1:proof:item:a}
\textit{a)} The equivalence of \SGARefLink{sga3:xiii:lemma:1:1:item:i}{(i)} and \SGARefLink{sga3:xiii:lemma:1:1:item:i-bis}{(i bis)} follows immediately
from \SGARefLink{sga3:xiii:equation:1}{(1)} and from the fact noted above that \(N\) is smooth over \(k\)
if and only if
\(\dim N=\operatorname{rank}_k\mathfrak n\). Consider the inclusion
morphism \(N\to M_a\). It is well known\footnote{Editor's note in the
French source: see, for example, EGA IV\textsubscript{4},
Theorem 17.11.1(d).} that if \(N\) is smooth over \(k\) at \(e\) and
the tangent map at \(e\) is surjective, then \(N\to M_a\) is smooth
at \(e\); since it is an immersion, it is therefore an isomorphism at
\(e\). This proves that \SGARefLink{sga3:xiii:lemma:1:1:item:i}{(i)} implies \SGARefLink{sga3:xiii:lemma:1:1:item:iii}{(iii)}.

To prove the equivalence of \SGARefLink{sga3:xiii:lemma:1:1:item:ii}{(ii)} and \SGARefLink{sga3:xiii:lemma:1:1:item:iii}{(iii)}, let
\(X_a=\psi^{-1}(a)\), as above. From the isomorphisms
\[
  M_a\simeq\varphi^{-1}(a)\simeq q^{-1}(X_a)
\]
we obtain a morphism \(p_a:M_a\to X_a\) making \(M_a\) a principal
homogeneous bundle with group \(N_{X_a}\). Consider \SGARefLink{sga3:xiii:lemma:1:1:diagram:01}{the diagram}
\SGARefTarget{sga3:xiii:lemma:1:1:diagram:01}
\[
\begin{tikzcd}[column sep=large,row sep=large]
M_a \arrow[d,"p_a"'] &
N \arrow[l,"j_a"'] \arrow[d] \\
X_a \arrow[r,"j'_a"] &
S=\operatorname{Spec}(k).
\end{tikzcd}
\]
Here \(\operatorname{Spec}(k)\to X_a\)\footnote{\SGARefLink{sga3:xiii:lemma:1:1:diagram:01}{The printed diagram}
directs \(j'_a:X_a\to S\), whereas the immediately following prose
defines the point morphism \(S\to X_a\) and calls it \(j'_a\). Both
\SGARefLink{sga3:xiii:lemma:1:1:diagram:01}{the printed diagram direction} and the verbal definition are retained
here. The source also prints \(X_n\) here and \(X_{a,b}\) in the two
local rings below. They are corrected to \(X_a\) and
\(X_{a,\bar a}\), in accordance with the fiber notation and the point
\(\bar a\).} is defined by the point
\(\bar a=(\bar e,a)\) of \(X\), and \(j_a:N\to M_a\) is the canonical
immersion. To say that \(\psi\) is unramified at \((\bar e,a)\) means
that \(j'_a\) is an open immersion, or equivalently that it induces an
isomorphism
\SGARefTarget{sga3:xiii:lemma:1:1:display:residue-field-map}
\[
  S\xrightarrow{\sim}\operatorname{Spec}
  \bigl(\mathcal O_{X_a,\bar a}\bigr).
\]
Since \(p_a\) is flat, this is equivalent to saying that the morphism
obtained from \SGARefLink{sga3:xiii:lemma:1:1:display:residue-field-map}{the preceding one} by the base change
\[
  \operatorname{Spec}\bigl(\mathcal O_{M_a,e}\bigr)
  \longrightarrow
  \operatorname{Spec}\bigl(\mathcal O_{X_a,\bar a}\bigr)
\]
is an isomorphism. The resulting morphism is precisely
\[
  \operatorname{Spec}\bigl(\mathcal O_{N,e}\bigr)
  \longrightarrow
  \operatorname{Spec}\bigl(\mathcal O_{M_a,e}\bigr),
\]
which proves the equivalence of \SGARefLink{sga3:xiii:lemma:1:1:item:ii}{(ii)} and \SGARefLink{sga3:xiii:lemma:1:1:item:iii}{(iii)}.

Notice moreover that the proof shows that \SGARefLink{sga3:xiii:lemma:1:1:item:ii}{(ii)} and \SGARefLink{sga3:xiii:lemma:1:1:item:iii}{(iii)} imply \SGARefLink{sga3:xiii:lemma:1:1:item:iii-bis}{the
following condition}, apparently stronger than \SGARefLink{sga3:xiii:lemma:1:1:item:iii}{(iii)}:
\begin{enumerate}
\item[(iii bis)] \SGARefTarget{sga3:xiii:lemma:1:1:item:iii-bis}
  \(N\) is an open and closed subscheme of \(M_a\).
\end{enumerate}

\SGARefTarget{sga3:xiii:lemma:1:1:proof:item:b}
\textit{b)} The first assertion follows from the isomorphism
\(M_a\simeq\varphi^{-1}(a)\), and the second from the fact that \(q\)
is flat and \(q(e,a)=(\bar e,a)\).

\section{Density theorem and theory of the regular points of \(G\)}
\SGARefTarget{sga3:xiii:section:2}

We apply the constructions and notation of \SGARefLink{sga3:xiii:section:1}{the preceding section} in
the following situation. Let \(G\) be a smooth connected algebraic
group over \(k\); take \(V=G\), with \(G\) acting on \(V\) by inner
automorphisms; and let \(W=H\), where \(H\) is a smooth connected
algebraic subgroup of \(G\). We denote the Lie algebras of \(G\) and
\(H\) by \(\mathfrak g\) and \(\mathfrak h\), the normalizer of \(H\)
in \(G\) by \(N\), and its Lie algebra by \(\mathfrak n\). For
\(a\in G(k)\), we again denote by \(M_a\), as in \SGARefLink{sga3:xiii:section:1}{Section 1}, the
inverse of the transporter of \(a\) into \(H\). Thus, if
\(a\in H(k)\), then \(N\subset M_a\); in that case
\(\mathfrak m_a\) denotes the Zariski tangent space of \(M_a\) at the
identity element \(e\) of \(G\). Note that
\[
  \mathfrak h\subset\mathfrak n\subset\mathfrak m_a
  \subset\mathfrak g
  \qquad\text{for }a\in H(k).
\]
We shall use \SGARefLink{sga3:xiii:lemma:2:0}{the following lemma}.

\medskip
\noindent\SGARefTarget{sga3:xiii:lemma:2:0}\textbf{Lemma 2.0.} ---
\textit{The equality \(H=N^0\) holds if and only if}
\SGARefTarget{sga3:xiii:lemma:2:0:display:invariants-zero}
\[
  (\mathfrak g/\mathfrak h)^H=0,
\]
\textit{where the left-hand side denotes the subspace of invariants
for the action of \(H\) induced by the adjoint representation. When
\SGARefLink{sga3:xiii:lemma:2:0:display:invariants-zero}{this condition} holds, \(N\) is smooth and \(\dim X=\dim G\). In all
cases, \(\dim X\leq\dim G\), and equality holds if and only if \(H\)
has finite index in \(N\).}

Indeed, by \SGARefLink{sga3:ii:theorem:5-2-3:item:i}{II~5.2.3(i)}, \(\mathfrak n\) is the inverse image of
\((\mathfrak g/\mathfrak h)^H\) under
\(\mathfrak g\to\mathfrak g/\mathfrak h\). Hence
\((\mathfrak g/\mathfrak h)^H=0\) is equivalent to
\(\mathfrak h=\mathfrak n\), which in turn is equivalent to
\(H=N^0\), since \(H\) is a smooth connected algebraic subgroup of
the algebraic group \(N\) (cf. VI.2). This also shows that \(N\) is
smooth. Furthermore,
\[
  \dim X=(\dim G-\dim N)+\dim H
        =\dim G-(\dim N-\dim H),
\]
and therefore
\[
  \dim X\leq\dim G,
\]
with equality if and only if \(\dim H=\dim N\), that is, if and only
if \(H\) has finite index in \(N\). This holds in particular when
\(H=N^0\), and proves the lemma.

\medskip
\noindent\SGARefTarget{sga3:xiii:theorem:2:1}\textbf{Theorem 2.1.} ---
\textit{Let \(G\) be a smooth connected algebraic group over an
algebraically closed field \(k\), let \(H\) be a smooth connected
algebraic subgroup, and let \(N\) be its normalizer. Write
\(\mathfrak g,\mathfrak h,\mathfrak n\) for their Lie algebras, let
\(X=G\times^N H\) be the scheme introduced in \SGARefLink{sga3:xiii:section:1}{Section 1} (fibered over
\(G/N\) with typical fiber \(H\)), and let \(\psi:X\to G\) be the
canonical morphism. Its image is also the image of
\(\varphi:G\times H\to G\), where}
\[
  \varphi(g,h)=\operatorname{int}(g)h=ghg^{-1}.
\]
\textit{The following conditions are equivalent:}

\begin{enumerate}
\item[(i)] \SGARefTarget{sga3:xiii:theorem:2:1:item:i}
  \(H\) contains a Cartan subgroup \(C\) of \(G\) (XII~1).
\item[(i bis)] \SGARefTarget{sga3:xiii:theorem:2:1:item:i-bis}
  \(H\) has the same reductive rank and the same nilpotent rank
  (XII~1) as \(G\).
\item[(ii)] \SGARefTarget{sga3:xiii:theorem:2:1:item:ii}
  \(H\) contains a maximal torus \(T\) of \(G\), and
  \((\mathfrak g/\mathfrak h)^T=0\).
\item[(iii)] \SGARefTarget{sga3:xiii:theorem:2:1:item:iii}
  The set of conjugates of \(H\) containing a given maximal torus is
  finite and nonempty, and \(H\) has finite index in its normalizer.
\item[(iv)] \SGARefTarget{sga3:xiii:theorem:2:1:item:iv}
  There exists \(a\in H(k)\) that belongs to only finitely many
  conjugates of \(H\) (or merely such that \(\psi^{-1}(a)\) has an
  isolated point), and \(H\) has finite index in its normalizer.
\item[(iv bis)] \SGARefTarget{sga3:xiii:theorem:2:1:item:iv-bis}
  The morphism \(\psi:X\to G\) is generically quasi-finite (that is,
  there is a dense open subset of \(X\) on which \(\psi\) is
  quasi-finite), and \(H\) has finite index in its normalizer.
\item[(v)] \SGARefTarget{sga3:xiii:theorem:2:1:item:v}
  There is a dense open subset \(U\) of \(G\) such that, for every
  \(x\in U(k)\), the set of conjugates of \(H\) containing \(x\) is
  finite and nonempty; equivalently, \(\psi:X\to G\) is dominant and
  generically quasi-finite.
\item[(vi)] \SGARefTarget{sga3:xiii:theorem:2:1:item:vi}
  There is a dense open subset \(U\) of \(G\) such that every
  \(x\in U(k)\) is contained in a conjugate of \(H\); equivalently,
  \(\psi:X\to G\) is dominant.
\item[(vii)] \SGARefTarget{sga3:xiii:theorem:2:1:item:vii}
  There exists \(a\in H(k)\) such that the fixed subspace of
  \(\mathfrak g/\mathfrak h\) under
  \(\operatorname{ad}_{\mathfrak g/\mathfrak h}(a)\) is zero.
\end{enumerate}

\textit{Moreover, these conditions imply that \(H\) is its own
connected normalizer: \(N\) is smooth and
\(\dim H=\dim N\). They also imply that \(\psi:X\to G\) is
generically étale.}

\noindent\textit{Proof.}
By \SGARefLink{sga3:xiii:lemma:2:0}{2.0}, \(\dim X\leq\dim G\), with equality if and only if
\(\dim H=\dim N\), or equivalently if and only if \(H\) has finite
index in \(N\). The inequality \(\dim X\leq\dim G\) shows that
\(\psi\) is dominant if and only if it is dominant and generically
quasi-finite, and also if and only if it is generically quasi-finite
and \(\dim X=\dim G\). Since the last equality is equivalent, by
\SGARefLink{sga3:xiii:lemma:2:0}{2.0}, to \(H\) having finite index in \(N\), conditions \SGARefLink{sga3:xiii:theorem:2:1:item:vi}{(vi)}, \SGARefLink{sga3:xiii:theorem:2:1:item:v}{(v)}, and
\SGARefLink{sga3:xiii:theorem:2:1:item:iv-bis}{(iv bis)} are equivalent. The equivalence of \SGARefLink{sga3:xiii:theorem:2:1:item:iv}{(iv)} and \SGARefLink{sga3:xiii:theorem:2:1:item:iv-bis}{(iv bis)} is
immediate.

The equivalence of \SGARefLink{sga3:xiii:theorem:2:1:item:i}{(i)} and \SGARefLink{sga3:xiii:theorem:2:1:item:i-bis}{(i bis)} follows directly from the
definitions and is left to the reader. If \(H\) contains a Cartan
subgroup \(C\) of \(G\), it contains the maximal torus \(T\) of
\(C\), which is a maximal torus of \(G\). Since \(C\) is the
centralizer of \(T\), its Lie algebra \(\mathfrak c\) is
\[
  \mathfrak c=\mathfrak g^T
\]
(\SGARefLink{sga3:ii:theorem:5-2-3:item:ii}{II~5.2.3(ii)}). Hence
\(\mathfrak g^T=\mathfrak c\subset\mathfrak h\), and \SGARefLink{sga3:i:proposition:4-7-3}{I~4.7.3} gives
the relation\footnote{The French source has two evident typographical
slips in this sentence, writing \(\mathfrak n\) for \(\mathfrak c\)
and reversing the inclusion involving \(\mathfrak h\) and
\(\mathfrak g^T\). The mathematically required relation is restored
here.}
\SGARefTarget{sga3:xiii:theorem:2:1:equation:x}
\[
\tag{x}
  (\mathfrak g/\mathfrak h)^T=0.
\]
Conversely, suppose that \(H\) contains the maximal torus \(T\) and
that \SGARefLink{sga3:xiii:theorem:2:1:equation:x}{(x)} holds, so that \(\mathfrak h\supset\mathfrak c\). We claim
that \(H\) contains the centralizer \(C\) of \(T\), which will prove
the equivalence of \SGARefLink{sga3:xiii:theorem:2:1:item:i}{(i)} and \SGARefLink{sga3:xiii:theorem:2:1:item:ii}{(ii)}. This follows from \SGARefLink{sga3:xiii:lemma:2:1:1}{the next lemma}.

\medskip
\noindent\SGARefTarget{sga3:xiii:lemma:2:1:1}\textbf{Lemma 2.1.1.} ---
\textit{Let \(G\) be a smooth algebraic group over a field \(k\), let
\(T\) be a subgroup of multiplicative type of \(G\), let \(C\) be its
connected centralizer (equal to the centralizer of \(T\) when \(G\)
is connected and \(T\) is a torus, by \SGARefLink{sga3:xii:6.6:item:b}{XII~6.6(b)}), and let \(H\) be
a smooth subgroup of \(G\) containing \(T\). Then \(H\) contains
\(C\) if and only if its Lie algebra \(\mathfrak h\) contains the Lie
algebra \(\mathfrak c\) of \(C\).}

Indeed, \(\operatorname{Centr}_G(T)\) is smooth over \(k\) by \SGARefLink{sga3:xi:corollary:2:4}{XI~2.4},
and hence \(C\) is smooth over \(k\). Likewise,
\(\operatorname{Centr}_H(T)\) is smooth, and
\[
  \operatorname{Centr}_H(T)=\operatorname{Centr}_G(T)\cap H
\]
has Lie algebra \(\mathfrak h\cap\mathfrak c\). The hypothesis
therefore implies that this smooth subgroup has the same Lie algebra
as the smooth group \(\operatorname{Centr}_G(T)\), so it contains the
connected component \(C\) of the latter. Thus \(H\) contains \(C\).
\hfill\(\square\)

We now prove the equivalence of \SGARefLink{sga3:xiii:theorem:2:1:item:i}{(i)} and \SGARefLink{sga3:xiii:theorem:2:1:item:iii}{(iii)}. If \(H\) contains the
maximal torus \(T\) of \(G\), we must show that \(H\supset C\)
(equivalently, that \SGARefLink{sga3:xiii:theorem:2:1:equation:x}{(x)} holds) if and only if \(H\) has finite index
in its normalizer and the set of conjugates of \(H\) containing \(T\)
is finite. If \(H\supset C\), then \SGARefLink{sga3:xiii:theorem:2:1:equation:x}{(x)} implies, a fortiori,
\SGARefTarget{sga3:xiii:theorem:2:1:equation:xx}
\[
\tag{xx}
  (\mathfrak g/\mathfrak h)^H=0.
\]
The Lie algebra \(\mathfrak n\) is the inverse image of the left-hand
side under the canonical homomorphism
\(\mathfrak g\to\mathfrak g/\mathfrak h\) (II~5.2.3(i)).
Consequently, by \SGARefLink{sga3:xiii:lemma:2:0}{2.0}, \SGARefLink{sga3:xiii:theorem:2:1:equation:xx}{(xx)} says that \(H=N^0\); in particular, \(H\)
has finite index in its normalizer. Consider \SGARefLink{sga3:xiii:theorem:2:1:diagram:01}{the diagram of
subgroups}
\SGARefTarget{sga3:xiii:theorem:2:1:diagram:01}
\[
\begin{tikzcd}[column sep=large,row sep=large]
T \arrow[r] &
N(T)\cap H \arrow[r] \arrow[d] &
H \arrow[d] \\
{} &
N(T)\cap N(H) \arrow[r] \arrow[d] &
N(H) \\
{} &
N(T). &
\end{tikzcd}
\]
By the conjugacy theorem for maximal tori in \(H\) (\SGARefLink{sga3:xii:6.6:item:a}{XII~6.6(a)}),
every conjugate of \(H\) containing \(T\) is obtained from \(H\) by
conjugation by an element of \(N(T)(k)\). Thus the set of such
conjugates is in bijection with
\[
  \bigl(N(T)/(N(T)\cap N(H))\bigr)(k).
\]
Since \(H\supset C\), we have \(N(T)\cap H\supset C\), so the
preceding set is a quotient of \((N(T)/C)(k)\), which is finite.
This proves \SGARefLink{sga3:xiii:theorem:2:1:item:i}{(i)} \(\Rightarrow\) \SGARefLink{sga3:xiii:theorem:2:1:item:iii}{(iii)}.

Conversely, suppose \SGARefLink{sga3:xiii:theorem:2:1:item:iii}{(iii)}, so that
\(N(T)/(N(T)\cap N(H))\) and \(N(H)/H\) are finite. The conjugacy
theorem in \(H\) shows that the homomorphism
\[
  (N(T)\cap N(H))/(N(T)\cap H)\longrightarrow N(H)/H
\]
induced by \SGARefLink{sga3:xiii:theorem:2:1:diagram:01}{the preceding diagram} is bijective on \(k\)-valued points
(indeed, it is an isomorphism). Since the latter quotient is finite,
so is the former. Hence \(N(T)\cap H\) has finite index in \(N(T)\)
and therefore contains \(C=N(T)^0\). Thus \(H\supset C\), and \SGARefLink{sga3:xiii:theorem:2:1:item:i}{(i)},
\SGARefLink{sga3:xiii:theorem:2:1:item:i-bis}{(i bis)}, \SGARefLink{sga3:xiii:theorem:2:1:item:ii}{(ii)}, and \SGARefLink{sga3:xiii:theorem:2:1:item:iii}{(iii)} are equivalent.

We next prove \SGARefLink{sga3:xiii:theorem:2:1:item:ii}{(ii)} \(\Rightarrow\) \SGARefLink{sga3:xiii:theorem:2:1:item:vii}{(vii)}. Both conditions are
invariant under an extension \(k\to k'\) with \(k'\) algebraically
closed. We may therefore assume that \(k\) has infinite
transcendence degree over its prime subfield. There is then an
element \(a\in T(k)\) such that the subgroup of \(T(k)\) generated by
\(a\) is Zariski dense in \(T\). It follows that
\[
  (\mathfrak g/\mathfrak h)^T
  =(\mathfrak g/\mathfrak h)^{\operatorname{ad}(a)}.
\]
The left-hand side is zero by hypothesis, which proves \SGARefLink{sga3:xiii:theorem:2:1:item:vii}{(vii)}.

The implication \SGARefLink{sga3:xiii:theorem:2:1:item:vii}{(vii)} \(\Rightarrow\) \SGARefLink{sga3:xiii:theorem:2:1:item:vi}{(vi)} is contained in \SGARefLink{sga3:xiii:corollary:2:2}{the
following result}, which sharpens \SGARefLink{sga3:xiii:theorem:2:1}{2.1}.

\medskip
\noindent\SGARefTarget{sga3:xiii:corollary:2:2}\textbf{Corollary 2.2.} ---
\textit{Let \(G\) be a smooth algebraic group over a field \(k\), let
\(H\) be a smooth algebraic subgroup, let \(N\) be its normalizer in
\(G\), and let}
\[
  \varphi:G\times H\longrightarrow G,\qquad
  \varphi(g,h)=\operatorname{ad}(g)h=ghg^{-1}.
\]
\textit{Let \(\psi:X=G\times^N H\to G\) be the morphism obtained from
\(\varphi\) by passage to the quotient (cf. \SGARefLink{sga3:xiii:section:1}{Section 1}), and let
\(a\in H(k)\). The following conditions are equivalent:}

\begin{enumerate}
\item[(i)] \SGARefTarget{sga3:xiii:corollary:2:2:item:i}
  \(\varphi\) is smooth at \((e,a)\).
\item[(ii)] \SGARefTarget{sga3:xiii:corollary:2:2:item:ii}
  \(\psi\) is étale at \((\bar e,a)\), and \(N\) is smooth over \(k\).
\item[(iii)] \SGARefTarget{sga3:xiii:corollary:2:2:item:iii}
  \((\mathfrak g/\mathfrak h)^{\operatorname{ad}(a)}=0\), where
  \(\mathfrak g\) and \(\mathfrak h\) are the Lie algebras of \(G\)
  and \(H\).
\end{enumerate}

\textit{These conditions imply \(H^0=N^0\).}

The smoothness at a \(k\)-rational point of \(\varphi\), a morphism
between smooth \(k\)-preschemes, is equivalent to the surjectivity of
the tangent map at that point. Under the usual identifications of
the tangent spaces of \(G\) and \(H\) with their Lie algebras, a
direct calculation gives
\SGARefTarget{sga3:xiii:corollary:2:2:display:tangent-map}
\[
  d\varphi(\xi,\eta)
  =\bigl(\operatorname{id}-\operatorname{ad}(a)\bigr)\xi+\eta,
\]
as a map from \(\mathfrak g\times\mathfrak h\) to \(\mathfrak g\).
Its surjectivity is therefore equivalent to the surjectivity of
\(\operatorname{id}-\operatorname{ad}(a)\) on
\(\mathfrak g/\mathfrak h\), which is precisely \SGARefLink{sga3:xiii:corollary:2:2:item:iii}{(iii)}.

Condition \SGARefLink{sga3:xiii:corollary:2:2:item:iii}{(iii)} implies, a fortiori,
\((\mathfrak g/\mathfrak h)^H=0\), and hence \(H^0=N^0\) by \SGARefLink{sga3:xiii:lemma:2:0}{2.0}
(the connectedness hypothesis made at the beginning of the section
is not needed there). It follows that \(N\) is smooth and that
\(\dim H=\dim N\), whence \(\dim X=\dim G\). Since
\(q:G\times H\to X\) is flat and
\(\psi\circ q=\varphi\)\footnote{The source reverses this composition,
printing \(\psi=\varphi\circ q\). The displayed morphisms and the
definition of \(\varphi\) require \(\varphi=\psi\circ q\).} is smooth
at \((e,a)\), the morphism \(\psi\)
is smooth at \(q(e,a)=(\bar e,a)\), and hence is étale there by
dimension. We have proved
\(\SGARefLink{sga3:xiii:corollary:2:2:item:i}{\mathrm{(i)}}\Longleftrightarrow\SGARefLink{sga3:xiii:corollary:2:2:item:iii}{\mathrm{(iii)}}
\Longrightarrow\SGARefLink{sga3:xiii:corollary:2:2:item:ii}{\mathrm{(ii)}}\). Conversely, \SGARefLink{sga3:xiii:corollary:2:2:item:ii}{(ii)} implies \SGARefLink{sga3:xiii:corollary:2:2:item:i}{(i)},
because the smoothness of \(N\) implies that of \(q\).
\hfill\(\square\)

In particular, under condition \SGARefLink{sga3:xiii:theorem:2:1:item:vii}{(vii)} of \SGARefLink{sga3:xiii:theorem:2:1}{2.1} the morphism
\(\varphi\) is smooth at \((e,a)\), so its image contains a nonempty
open subset of \(G\). Thus \(\psi\) is dominant, proving
\SGARefLink{sga3:xiii:theorem:2:1:item:vii}{(vii)} \(\Rightarrow\) \SGARefLink{sga3:xiii:theorem:2:1:item:vi}{(vi)}.

It remains to prove \SGARefLink{sga3:xiii:theorem:2:1:item:vi}{(vi)} \(\Rightarrow\) \SGARefLink{sga3:xiii:theorem:2:1:item:i}{(i)}. Suppose first that
\(G\) is affine. Let \(U\) be a nonempty open subset of \(G\) such
that every \(x\in U(k)\) is contained in a conjugate of \(H\), and
let \(C\) be a Cartan subgroup of \(G\). Applying the already proved
implication \SGARefLink{sga3:xiii:theorem:2:1:item:i}{(i)} \(\Rightarrow\) \SGARefLink{sga3:xiii:theorem:2:1:item:iv}{(iv)} with \(C\) in place of \(H\)
(Borel's ``density theorem''), we find a conjugate of \(C\) meeting
\(U\). We may therefore assume that \(U\cap C\neq\varnothing\);
equivalently, there is a nonempty open subset \(V\) of \(C\) such
that every \(x\in V(k)\) belongs to a conjugate of \(H\).

Write
\SGARefTarget{sga3:xiii:theorem:2:1:display:cartan-decomposition}
\[
  C=T\cdot C_u,
\]
where \(T\) is the maximal torus of \(C\), hence a maximal torus of
\(G\), and \(C_u\) is the unipotent part of \(C\); \(T\) lies in the
center of \(C\) (\emph{BIBLE} 6, Theorem 2). We may again assume that
\(k\) has infinite transcendence degree over its prime subfield.
Choose \(t\in T(k)\) in the projection of \(V\) onto \(T\), a
nonempty open subset of \(T\), so that
\(tC_u\cap V\neq\varnothing\), and choose it so that it generates a
Zariski-dense subgroup of \(T\). Every algebraic subgroup of \(G\)
containing a product \(tu\), with \(t\in T(k)\) and
\(u\in C_u(k)\), contains both factors (\emph{BIBLE} 4,
Theorem 3). Taking \(tu\in V(k)\), we conclude that a conjugate of
\(H\) contains \(t\), and therefore contains \(T\). After replacing
\(H\) by that conjugate, we may assume
\[
  T\subset H.
\]

Let \(W\) be the inverse image in \(C_u\) of \(V\) under
\(u\mapsto tu\). For every \(x\in(T\cdot W)(k)\), there is then a conjugate
of \(H\) containing both \(T\) and \(x\). As observed above, every
such conjugate is of the form \(\operatorname{int}(g)H\) for some
\(g\in N(T)(k)\). Consider
\SGARefTarget{sga3:xiii:theorem:2:1:display:f}
\[
  f:N(T)\times H\longrightarrow G,\qquad
  f(g,h)=\operatorname{int}(g)h=ghg^{-1}.
\]
The image of \(f\) contains \(TW\). Since \(C\) has finite index in
\(N(T)\), the latter is a finite union of translates \(Cg_i\), with
\(g_i\in N(T)(k)\). It follows that some dense open subset \(V'\) of
\(C=T C_u\) is contained in the image under \(f\) of
\((Cg_i)\times H\). Replacing \(H\) by
\(\operatorname{int}(g_i)H\), we may assume \(g_i=e\), so that
\[
  f(C\times H)\supset V'.
\]
Thus for every \(u\in V'(k)\) there are \(v\in C(k)\) and
\(h\in H(k)\) such that
\[
  v^{-1}hv=u,\qquad\text{hence}\qquad vuv^{-1}\in H(k).
\]
Set
\[
  C'=C\cap H=\operatorname{Centr}_H(T).
\]
Since \(vuv^{-1}\) also belongs to \(C\), it lies in \(C'\), and
hence \(u\in\operatorname{int}(v^{-1})C'\). Therefore the union of
the conjugates of \(C'\) in \(C\), by elements of \(C(k)\), is
dense. As in the preceding argument for the pair \((G,H)\), now
applied to \((C,C')\), this implies that \(C'\) has finite index in
its normalizer in \(C\). By the lemma of \emph{BIBLE} 7, \SGARefLink{sga3:xiii:section:1}{Section 1},
we obtain \(C'=C\), and consequently \(C\subset H\).\footnote{The
source prints \(C=H\); only the inclusion \(C\subset H\) follows and
is used.} This proves
\SGARefLink{sga3:xiii:theorem:2:1:item:vi}{(vi)} \(\Rightarrow\) \SGARefLink{sga3:xiii:theorem:2:1:item:i}{(i)} when \(G\) is affine.

In the general case we argue by induction on \(n=\dim G\), the
assertion being trivial for \(n=0\). Let \(Z\) be the center of
\(G\). There are two cases.

\begin{enumerate}
\item[\(1^\circ\)] \SGARefTarget{sga3:xiii:theorem:2:1:proof:case:1} If \(\dim(Z\cap H)>0\), set
  \(G'=G/(Z\cap H)\). Then \(\dim G'<n\), and condition \SGARefLink{sga3:xiii:theorem:2:1:item:vi}{(vi)} for
  \(H\) implies \SGARefLink{sga3:xiii:theorem:2:1:item:vi}{the same condition} for its image \(H'\) in \(G'\).
  By induction, \(H'\) contains a Cartan subgroup \(C'\) of \(G'\).
  Hence \(H\) contains the inverse image \(C\) of \(C'\), which is a
  Cartan subgroup by \SGARefLink{sga3:xii:theorem:6-6:item:e}{XII~6.6(e)}.

\item[\(2^\circ\)] \SGARefTarget{sga3:xiii:theorem:2:1:proof:case:2} If \(\dim(Z\cap H)=0\), the canonical morphism
  \(H\to G/Z\) is finite. Since \(G/Z\) is affine by \SGARefLink{sga3:xii:6.1}{XII~6.1}, \(H\)
  is affine. Every homomorphism from \(H\) to an abelian variety is
  therefore zero (indeed, every morphism of preschemes from \(H\) to
  an abelian variety is constant). This follows from the fact that a
  smooth connected affine algebraic group over an algebraically
  closed field is a rational variety, or simply from the fact that
  it is the union of its Borel subgroups (\emph{BIBLE} 6,
  Theorem 5(b)); the structure theorems \emph{BIBLE} 6.2 and 6.3
  show readily that a smooth connected solvable affine group is a
  rational variety.

  Apply Chevalley's structure theorem to \(G\): the group \(G\) is
  an extension of an abelian variety \(A\) by a smooth affine group.
  Since \(H\) is affine, its image in \(A\) is zero. On the other
  hand, that image must be all of \(A\), because the union of its
  conjugates in \(A\) is dense. Hence \(A=0\), so \(G\) is affine,
  reducing us to the case already proved.
\end{enumerate}

This completes the proof of \SGARefLink{sga3:xiii:theorem:2:1}{2.1}.

\medskip
\noindent\SGARefTarget{sga3:xiii:corollary:2:3}\textbf{Corollary 2.3.} ---
\textit{Assume that the equivalent conditions of \SGARefLink{sga3:xiii:theorem:2:1}{2.1} hold.}

\begin{enumerate}
\item[\textit{a)}] \SGARefTarget{sga3:xiii:corollary:2:3:item:a}
  Let \(k(X)\) and \(k(G)\) be the fields of rational functions of
  \(X\) and \(G\), respectively. Then \(k(X)\) is a finite separable
  extension of \(k(G)\). Denote its degree by \(d\).

\item[\textit{b)}] \SGARefTarget{sga3:xiii:corollary:2:3:item:b}
  Let \(T\) be a maximal torus of \(G\) contained in \(H\), which
  exists by \SGARefLink{sga3:xiii:theorem:2:1:item:ii}{2.1(ii)}, and let \(C\) be the corresponding Cartan
  subgroup of \(G\). Then \(C\subset H\). Moreover,
  \(\operatorname{Norm}_G(T)\) is a smooth subgroup of \(G\), and
  \[
    \operatorname{Norm}_G(T)\cap\operatorname{Norm}_G(H)
    =
    \operatorname{Norm}_{\operatorname{Norm}_G(H)}(T)
  \]
  is a smooth subgroup of it, of finite index equal to the degree
  \(d\) defined in \SGARefLink{sga3:xiii:corollary:2:3:item:a}{(a)}. The number of conjugates of \(H\) containing
  a given maximal torus, or a given Cartan subgroup, is \(d\).

\item[\textit{c)}] \SGARefTarget{sga3:xiii:corollary:2:3:item:c}
  Let \(U\) be the largest open subset of \(G\) such that
  \(\psi:X\to G\) induces a finite étale morphism
  \[
    \psi^{-1}(U)\longrightarrow U.
  \]
  Then \(U\) is dense. For \(g\in G(k)\), the following are
  equivalent:
  \begin{enumerate}
  \item \SGARefTarget{sga3:xiii:corollary:2:3:item:c:condition:i} \(g\in U(k)\);
  \item \SGARefTarget{sga3:xiii:corollary:2:3:item:c:condition:ii} exactly \(d\) conjugates of \(H\) contain \(g\);
  \item \SGARefTarget{sga3:xiii:corollary:2:3:item:c:condition:iii} there are at least \(d\) distinct conjugates \(H_i\) of \(H\)
    containing \(g\), and for every \(i\),
    \[
      (\mathfrak g/\mathfrak h_i)^{\operatorname{ad}(g)}=0,
      \qquad \mathfrak h_i=\operatorname{Lie}(H_i).
    \]
  \end{enumerate}
\end{enumerate}

Assertion \SGARefLink{sga3:xiii:corollary:2:3:item:a}{(a)} follows from the generic étaleness of \(\psi\), stated
at the end of \SGARefLink{sga3:xiii:theorem:2:1}{2.1}. This also shows that the open set \(U\) in \SGARefLink{sga3:xiii:corollary:2:3:item:c}{(c)} is
nonempty, hence dense, and gives the two stated characterizations of
its \(k\)-points: here one uses that \(\psi\) is separated, that
\(X\) is integral, that \(G\) is integral and normal (SGA 1,
I~10.11), and that
\((\mathfrak g/\mathfrak h_i)^{\operatorname{ad}(g)}=0\) means that
\(\psi\) is étale at the point \(x_i\in\psi^{-1}(g)\) corresponding
to \(H_i\).

If \(H\) contains the maximal torus \(T\) of \(G\), the centralizers
of \(T\) in \(H\) and \(G\) have the same dimension and are smooth
and connected (\SGARefLink{sga3:xii:theorem:6-6:item:b}{XII~6.6(b)}); they are therefore equal, proving that
\(C\subset H\).\footnote{The source prints the tautological
\(C\subset G\); equality of the two centralizers gives
\(C\subset H\).} The normalizer of \(T\) in any smooth group
containing it is smooth (\SGARefLink{sga3:xi:corollary:2:4:bis}{XI~2.4 bis}), so
\(\operatorname{Norm}_G(T)\) and
\(\operatorname{Norm}_{\operatorname{Norm}_G(H)}(T)\) are smooth.
Recall also that \(N=\operatorname{Norm}_G(H)\) is smooth by the
last assertion of \SGARefLink{sga3:xiii:theorem:2:1}{2.1}. Furthermore,
\(\operatorname{Norm}_N(T)\) contains \(C\), which has finite index
in \(\operatorname{Norm}_G(T)\); hence
\(\operatorname{Norm}_N(T)\) itself has finite index in
\(\operatorname{Norm}_G(T)\).

The conjugacy theorem for maximal tori in \(H\) shows that this index
is the number of conjugates of \(H\) containing \(T\), equivalently
the number containing \(C\). The union of the conjugates of \(C\) in
\(G\) is dense, by \SGARefLink{sga3:xiii:theorem:2:1:item:i}{2.1(i)} \(\Rightarrow\) \SGARefLink{sga3:xiii:theorem:2:1:item:vi}{(vi)} applied to \(C\) in
place of \(H\), and \(U\) is stable under inner automorphisms. Hence
\(C\cap U\neq\varnothing\). Arguing as in the proof of
\SGARefLink{sga3:xiii:theorem:2:1:item:vi}{2.1(vi)} \(\Rightarrow\) \SGARefLink{sga3:xiii:theorem:2:1:item:ii}{(ii)}, and making a harmless extension of the
base field if necessary, one finds
\(g\in(C\cap U)(k)\) such that every conjugate of \(H\) containing
\(g\) contains \(T\), and hence \(C\). Thus the conjugates of \(H\)
containing \(C\) are precisely those containing \(g\); since
\(g\in U(k)\), their number is \(d\). This proves \SGARefLink{sga3:xiii:corollary:2:3:item:b}{(b)}.

In fact, the argument shows that the set of conjugates of \(H\)
containing \(T\) is a homogeneous set under the group of rational
points of
\[
  W_G(T)=
  \operatorname{Norm}_G(T)/\operatorname{Centr}_G(T).
\]
In particular,
\[
  d\leq\text{the order of the Weyl group of \(G\)}.
\]

\medskip
\noindent\SGARefTarget{sga3:xiii:corollary:2:4}\textbf{Corollary 2.4.} ---
\textit{With the notation of \SGARefLink{sga3:xiii:theorem:2:1}{2.1}, the following conditions are
equivalent:}

\begin{enumerate}
\item[(i)] \SGARefTarget{sga3:xiii:corollary:2:4:item:i}
  \(\psi:X\to G\) is birational.
\item[(ii)] \SGARefTarget{sga3:xiii:corollary:2:4:item:ii}
  Exactly one conjugate of \(H\) contains a given Cartan subgroup of
  \(G\).
\item[(iii)] \SGARefTarget{sga3:xiii:corollary:2:4:item:iii}
  \(H\) contains a Cartan subgroup \(C\) of \(G\), and
  \(\operatorname{Norm}_G(H)\supset\operatorname{Norm}_G(C)\).
\item[(iv)] \SGARefTarget{sga3:xiii:corollary:2:4:item:iv}
  There is a nonempty open subset \(V\) of \(G\) such that every
  \(g\in V(k)\) belongs to exactly one conjugate of \(H\).
\end{enumerate}

This follows immediately from \SGARefLink{sga3:xiii:theorem:2:1}{2.1} and \SGARefLink{sga3:xiii:corollary:2:3}{2.3}.

\medskip
\noindent\SGARefTarget{sga3:xiii:corollary:2:5}\textbf{Corollary 2.5.} ---
\textit{Assume that the conditions of \SGARefLink{sga3:xiii:corollary:2:4}{2.4} hold, and let
\(g\in G(k)\). The following conditions are equivalent:}

\begin{enumerate}
\item[(i)] \SGARefTarget{sga3:xiii:corollary:2:5:item:i}
  \(g\in U(k)\), where \(U\) is defined in \SGARefLink{sga3:xiii:corollary:2:3:item:c}{2.3(c)}; that is, \(g\) is
  contained in exactly one conjugate of \(H\).
\item[(ii)] \SGARefTarget{sga3:xiii:corollary:2:5:item:ii}
  The set of conjugates of \(H\) containing \(g\) is finite and
  nonempty.
\item[(iii)] \SGARefTarget{sga3:xiii:corollary:2:5:item:iii}
  The scheme \(\psi^{-1}(g)\), the ``scheme of conjugates of \(H\)
  containing \(g\),'' has an isolated point.
\item[(iv)] \SGARefTarget{sga3:xiii:corollary:2:5:item:iv}
  There is a conjugate \(H'\) of \(H\) containing \(g\) such that
  \[
    (\mathfrak g/\mathfrak h')^{\operatorname{ad}(g)}=0,
    \qquad \mathfrak h'=\operatorname{Lie}(H').
  \]
\end{enumerate}

\textit{Finally, \(U\) is also the largest open subset of \(G\) on
which \(\psi\) induces an isomorphism}
\[
  \psi^{-1}(U)\xrightarrow{\sim}U.
\]

The equivalence of \SGARefLink{sga3:xiii:corollary:2:5:item:i}{(i)}, \SGARefLink{sga3:xiii:corollary:2:5:item:ii}{(ii)}, and \SGARefLink{sga3:xiii:corollary:2:5:item:iii}{(iii)}, together with the last
assertion, follows from the ``Main Theorem''\footnote{Editor's note in
the French source: Zariski's Main Theorem.} applied to the
birational morphism \(\psi:X\to G\), since \(G\) is normal. Their
equivalence with \SGARefLink{sga3:xiii:corollary:2:5:item:iv}{(iv)} follows immediately from \SGARefLink{sga3:xiii:corollary:2:2}{Corollary~2.2}, which
characterizes the points of \(X\) at which \(\psi\) is
étale.\footnote{The source cites the final assertion of \SGARefLink{sga3:xiii:theorem:2:1}{2.1} here;
the pointwise étaleness criterion is \SGARefLink{sga3:xiii:corollary:2:2}{Corollary 2.2}.}

\medskip
\noindent\SGARefTarget{sga3:xiii:theorem:2:6}\textbf{Theorem 2.6.} ---
\textit{Let \(G\) be a smooth connected algebraic group over an
algebraically closed field \(k\), let \(C\) be a Cartan subgroup
associated with a maximal torus \(T\), and put}
\[
  N=\operatorname{Norm}_G(C)=\operatorname{Norm}_G(T)
  \qquad\text{(cf. XII~8.4).}
\]
\textit{Let \(X=G\times^N C\), where \(N\) acts on the left factor
\(G\) by right translations and on the right factor \(C\) by inner
automorphisms, and let \(\psi:X\to G\) be the canonical morphism.}

\begin{enumerate}
\item[\textit{a)}] \SGARefTarget{sga3:xiii:theorem:2:6:item:a}
  \(\psi\) is birational.

\item[\textit{b)}] \SGARefTarget{sga3:xiii:theorem:2:6:item:b}
  Let \(U\) be the largest open subset of \(G\) such that \(\psi\)
  induces an isomorphism
  \[
    \psi^{-1}(U)\xrightarrow{\sim}U
  \]
  (cf. \SGARefLink{sga3:xiii:corollary:2:5}{2.5}), and let
  \SGARefTarget{sga3:xiii:theorem:2:6:display:rho}
  \[
    \rho=\rho_\nu(G)=\dim C
  \]
  be the nilpotent rank of \(G\). For every \(g\in G(k)\), the
  multiplicity of the eigenvalue \(1\) in the action of
  \(\operatorname{ad}(g)\) on \(\mathfrak g\) is at least \(\rho\);
  it equals \(\rho\) if and only if \(g\in U(k)\).
\end{enumerate}

\noindent\textit{Proof.}
\SGARefLink{sga3:xiii:theorem:2:1:item:i}{Condition 2.1(i)} holds, so \SGARefLink{sga3:xiii:corollary:2:4:item:iii}{2.4(iii)} \(\Rightarrow\) \SGARefLink{sga3:xiii:corollary:2:4:item:i}{(i)} proves \SGARefLink{sga3:xiii:theorem:2:6:item:a}{(a)}.
In \emph{BIBLE} 7, when \(G\) is affine, the points of \(U(k)\) are
called the \emph{regular points} of \(G(k)\). We shall use the same
terminology and call \(U\) the \emph{open subset of regular points}
of \(G\). Notice that the proof in \emph{BIBLE} that this set is
itself open is incorrect; \SGARefLink{sga3:xiii:expose}{the present Exposé} establishes its
openness under more general hypotheses.

To prove \SGARefLink{sga3:xiii:theorem:2:6:item:b}{(b)}, for each \(g\in G(k)\) introduce \SGARefLink{sga3:xiii:theorem:2:6:display:characteristic-polynomial}{the characteristic
polynomial}
\SGARefTarget{sga3:xiii:theorem:2:6:display:characteristic-polynomial}
\[
  P(\operatorname{ad}(g),t)
  =t^n+c_1(g)t^{n-1}+\cdots+c_n(g).
\]
Replacing \(k\) by an arbitrary \(k\)-algebra shows immediately that
the functions \(c_i(g)\) arise from uniquely determined sections
\[
  c_i\in\Gamma(G,\mathcal O_G).
\]

Suppose that \(g\in G(k)\) belongs to a Cartan subgroup, which we may
take to be \(C\); this includes the case of a regular element. By
\SGARefLink{sga3:xiii:corollary:2:5:item:iv}{2.5(iv)},
\[
  (\mathfrak g/\mathfrak c)^{\operatorname{ad}(g)}=0
\]
if and only if \(g\) is regular, where
\(\mathfrak c=\operatorname{Lie}(C)\). Since \(C\) is nilpotent,
\(\operatorname{ad}_{\mathfrak c}(g)\) has only the eigenvalue \(1\).
It follows that the multiplicity of the eigenvalue \(1\) in
\(\operatorname{ad}_{\mathfrak g}(g)\) is at least \(\rho\), and is
exactly \(\dim C=\rho\) if and only if \(g\) is regular.

In particular, \SGARefLink{sga3:xiii:theorem:2:6:display:characteristic-polynomial}{the characteristic polynomial above} is divisible by
\((t-1)^\rho\). Divisibility by \((t-1)^\rho\) is expressed by
linear relations, with integer coefficients, among the coefficients
of the polynomial. These relations hold for \(g\in U(k)\), where
\(U\) is dense; since \(G\) is reduced, they therefore hold
everywhere. Thus\footnote{The French source prints \(c_0\) in place
of the omitted tail \(+\cdots+c_n\) on the left-hand side of
\SGARefLink{sga3:xiii:theorem:2:6:equation:dagger}{\((\dagger)\)}. \SGARefLink{sga3:xiii:theorem:2:6:display:characteristic-polynomial}{The characteristic polynomial displayed immediately
above} determines the correction made here.}
\SGARefTarget{sga3:xiii:theorem:2:6:equation:dagger}
\[
\tag{\(\dagger\)}
  t^n+c_1t^{n-1}+\cdots+c_n
  =(t-1)^\rho
    \bigl(t^{n-\rho}+b_1t^{n-\rho-1}+\cdots+b_{n-\rho}\bigr)
\]
in the polynomial ring over \(\Gamma(G,\mathcal O_G)\). In
particular, for every field extension \(K/k\) and every
\(g\in G(K)\), the eigenvalue \(1\) of \(\operatorname{ad}(g)\) has
multiplicity at least \(\rho\).

We have already seen that equality holds when \(g\) is regular; let
us prove the converse. Suppose first that \(G\) is affine, and write
\[
  g=g_sg_u
\]
as the product of its semisimple and unipotent parts
(\emph{BIBLE} 4, Section 4). Then
\[
  \operatorname{ad}(g)
  =\operatorname{ad}(g_s)\operatorname{ad}(g_u)
\]
is the corresponding decomposition of \(\operatorname{ad}(g)\)
(loc. cit., corollary to Theorem 3). Hence
\(\operatorname{ad}(g)\) and \(\operatorname{ad}(g_s)\) have the same
eigenvalues, counted with multiplicity; in particular, the
eigenvalue \(1\) has the same multiplicity in both.

By \emph{BIBLE} 7, Theorem 2, Corollary 1, the element \(g\) is
regular if and only if \(g_s\) is regular. We may therefore assume
that \(g=g_s\) is semisimple. It is then contained in a maximal torus
by \emph{BIBLE} 6, Theorem 5(c), and hence in a Cartan subgroup, a
case already treated. This proves \SGARefLink{sga3:xiii:theorem:2:6:item:b}{(b)} when \(G\) is affine.

In general, let \(Z=\operatorname{Centr}(G)_{\mathrm{red}}\). By
\SGARefLink{sga3:xii:6.6:item:e}{XII~6.6(e)}, the Cartan subgroups of \(G\) are the inverse images of
those of \(G'=G/Z\). Thus \(g\) is regular in \(G\) if and only if
its image \(g'\) is regular in \(G'\). Since \(Z\) is smooth, the Lie
algebra of \(G'\) is
\(\mathfrak g'=\mathfrak g/\mathfrak z\), where
\(\mathfrak z=\operatorname{Lie}(Z)\), and
\(\operatorname{ad}(g')\) is the endomorphism induced by
\(\operatorname{ad}(g)\) on \(\mathfrak g/\mathfrak z\).
Consequently, the multiplicity of the eigenvalue \(1\) in
\(\operatorname{ad}(g)\) is
\[
  d=\dim Z
\]
plus its multiplicity in \(\operatorname{ad}(g')\). The former
therefore equals the nilpotent rank of \(G\) if and only if the latter
equals the nilpotent rank of \(G'\). We are reduced to \(G'\), which
is affine by \SGARefLink{sga3:xii:6.1}{XII~6.1}, and that case has already been proved. This
completes the proof of \SGARefLink{sga3:xiii:theorem:2:6}{2.6}.

\medskip
\noindent\SGARefTarget{sga3:xiii:corollary:2:7}\textbf{Corollary 2.7.} ---
\textit{With the notation of the preceding proof,\footnote{Editor's
note in the French source: cf. \SGARefLink{sga3:xiii:theorem:2:6:equation:dagger}{\((\dagger)\)} above.} let}
\[
  b=1+b_1+\cdots+b_{n-\rho}\in\Gamma(G,\mathcal O_G).
\]
\textit{Then the open subset of regular points of \(G\) is}
\[
  U=G_b,
\]
\textit{the set of points of \(G\) at which \(b\) is invertible. In
particular, \(U\) is an affine open subset when \(G\) is affine.}

\medskip
\noindent\SGARefTarget{sga3:xiii:corollary:2:8}\textbf{Corollary 2.8.} ---
\textit{Let \(H\) be a smooth connected algebraic subgroup of \(G\)
containing a Cartan subgroup of \(G\).}

\begin{enumerate}
\item[\textit{a)}] \SGARefTarget{sga3:xiii:corollary:2:8:item:a}
  Let \(C\) be an algebraic subgroup of \(H\). Then \(C\) is a
  Cartan subgroup of \(H\) if and only if it is a Cartan subgroup of
  \(G\).

\item[\textit{b)}] \SGARefTarget{sga3:xiii:corollary:2:8:item:b}
  Let \(g\in G(k)\), and let \(d\) be the integer introduced in \SGARefLink{sga3:xiii:corollary:2:3}{2.3}.
  Then \(g\) is a regular point of \(G\) if and only if there are
  exactly \(d\) conjugates \(H_i\) of \(H\) containing \(g\), and
  \(g\) is a regular element of every \(H_i\). Equivalently, there
  are at most \(d\) conjugates of \(H\) containing \(g\), and \(g\)
  is regular in one of them. When these conditions hold, if \(C\) is
  the unique Cartan subgroup of \(G\) containing \(g\), then the
  conjugates of \(H\) containing \(g\) are precisely the conjugates
  of \(H\) containing \(C\).

\item[\textit{c)}] \SGARefTarget{sga3:xiii:corollary:2:8:item:c}
  Let \(g\in H(k)\). Then \(g\) is regular in \(G\) if and only if it
  is regular in \(H\) and
  \[
    (\mathfrak g/\mathfrak h)^{\operatorname{ad}(g)}=0.
  \]
\end{enumerate}

\noindent\textit{Proof.}
\SGARefTarget{sga3:xiii:corollary:2:8:proof:item:a}
\textit{a)} Under either hypothesis on \(C\), its unique maximal
torus \(T\) is a maximal torus of both \(H\) and \(G\), since \(H\)
has the same reductive rank as \(G\). The smooth connected groups
\(\operatorname{Centr}_H(T)\subset\operatorname{Centr}_G(T)\) have
the same dimension, and therefore are equal. Hence \(C\) is equal to
one of these two centralizers if and only if it is equal to the
other, proving \SGARefLink{sga3:xiii:corollary:2:8:item:a}{(a)}.

\SGARefTarget{sga3:xiii:corollary:2:8:proof:item:b}
\textit{b)} Suppose first that \(g\) is regular in \(G\), and let
\(C\) be the unique Cartan subgroup of \(G\) containing \(g\). By
\SGARefLink{sga3:xiii:corollary:2:3:item:b}{2.3(b)}, exactly \(d\) conjugates \(H_i\) of \(H\) contain \(C\).
Since
\[
  (\mathfrak g/\mathfrak c)^{\operatorname{ad}(g)}=0,
\]
the endomorphism \(\operatorname{ad}(g)\) has no eigenvalue \(1\) on
\(\mathfrak g/\mathfrak c\), and a fortiori
\[
  (\mathfrak g/\mathfrak h_i)^{\operatorname{ad}(g)}=0.
\]
By \SGARefLink{sga3:xiii:corollary:2:3:item:c}{2.3(c)}, exactly \(d\) conjugates of \(H\) contain \(g\), namely
the \(H_i\). A Cartan subgroup of such an \(H_i\) containing \(g\) is
a Cartan subgroup of \(G\) by \SGARefLink{sga3:xiii:corollary:2:8:item:a}{(a)}, and hence is \(C\). Thus \(g\) is
regular in every \(H_i\).

Conversely, suppose that at most \(d\) conjugates \(H_i\) of \(H\)
contain \(g\), and that \(g\) is regular in one of them, which we may
take to be \(H\). Since \(g\) is regular in \(H\), it belongs to a
unique Cartan subgroup \(C\) of \(H\); by \SGARefLink{sga3:xiii:corollary:2:8:item:a}{(a)}, \(C\) is a Cartan
subgroup of \(G\). Let \(C'\) be any Cartan subgroup of \(G\)
containing \(g\). By \SGARefLink{sga3:xiii:corollary:2:3:item:b}{2.3(b)}, exactly \(d\) conjugates of \(H\)
contain \(C'\). They all contain \(g\), so they are precisely the
\(H_i\); in particular, \(H\) contains \(C'\). Thus \(C\) and \(C'\)
are two Cartan subgroups of \(H\) containing the same regular element
\(g\) of \(H\), and therefore \(C=C'\). Hence \(g\) is regular in
\(G\).\footnote{The source prints ``regular in \(C\)'' at this
conclusion; the argument proves regularity in \(G\).}
\hfill\(\square\)

\SGARefTarget{sga3:xiii:corollary:2:8:proof:item:c}
\textit{c)} For an endomorphism \(u\) of a finite-dimensional vector
space, let \(\nu(u)\) denote the nullity\footnote{Editor's note in the
French source: that is, the dimension of the nullspace of
\(\operatorname{id}-u\).} of \(\operatorname{id}-u\). Then
\SGARefTarget{sga3:xiii:corollary:2:8:display:nullity}
\[
  \nu\bigl(\operatorname{ad}(g)_{\mathfrak g}\bigr)
  =
  \nu\bigl(\operatorname{ad}(g)_{\mathfrak h}\bigr)
  +
  \nu\bigl(\operatorname{ad}(g)_{\mathfrak g/\mathfrak h}\bigr).
\]
The two terms on the right are respectively at least the nilpotent
rank of \(H\), which equals the nilpotent rank \(\rho\) of \(G\) by
\SGARefLink{sga3:xiii:corollary:2:8:item:a}{(a)}, and at least \(0\). Hence
\(\nu(\operatorname{ad}(g)_{\mathfrak g})=\rho\) if and only if
\[
  \nu\bigl(\operatorname{ad}(g)_{\mathfrak h}\bigr)=\rho,
  \qquad
  \nu\bigl(\operatorname{ad}(g)_{\mathfrak g/\mathfrak h}\bigr)=0.
\]
This says precisely that \(g\) is regular in \(G\) if and only if it
is regular in \(H\) and
\(\operatorname{ad}(g)_{\mathfrak g/\mathfrak h}\) has no nonzero
invariants. \hfill\(\square\)

\medskip
\noindent\SGARefTarget{sga3:xiii:remark:2:9}\textbf{Remarks 2.9.} ---
In the statement of \SGARefLink{sga3:xiii:theorem:2:1}{2.1}, condition \SGARefLink{sga3:xiii:theorem:2:1:item:iii}{(iii)} cannot be weakened by
assuming only that \(H\) contains a maximal torus and has finite
index in its normalizer, even if one also requires that normalizer to
be smooth, that is, that \(H=N^0\), and even when \(G\) is affine and
solvable. An example is the group \(G\) of matrices of the form
\SGARefTarget{sga3:xiii:remark:2:9:display:matrix}
\[
  g=
  \begin{pmatrix}
    t & a & c\\
    0 & 1 & b\\
    0 & 0 & 1
  \end{pmatrix},
\]
and the subgroup \(H\) formed by \SGARefLink{sga3:xiii:remark:2:9:display:matrix}{the matrices above} for which
\(b=c=0\). The Cartan subgroup of \(G\) here consists of the matrices
\(g\) for which \(a=c=0\).

\medskip
\noindent\SGARefTarget{sga3:xiii:remark:2:10}\textbf{Remarks 2.10.} ---
Let \(G\) be a smooth algebraic group over \(k\), and let \(H\) be a
smooth algebraic subgroup; we no longer assume that \(G\) and \(H\)
are connected. Suppose that \(H^0\) contains a Cartan subgroup of
\(G^0\). Then
\[
  (\mathfrak g/\mathfrak h)^H
  \subset(\mathfrak g/\mathfrak h)^{H^0}=0,
\]
so \(H^0=N^0\), where \(N\) is the normalizer of \(H\); in
particular, \(N\) is smooth.

Nevertheless, even with \(G\) connected, examples are easily
constructed in which \(H\) has a connected component \(H_i\) such
that
\[
  (\mathfrak g/\mathfrak h)^{\operatorname{ad}(h)}\neq0
  \qquad\text{for every }h\in H_i(k).
\]
Equivalently, the morphism
\[
  G\times H_i\longrightarrow G,\qquad
  (g,h)\longmapsto\operatorname{ad}(g)\cdot h
\]
is nowhere étale, or even quasi-finite. One may take, for example,
\(H\) to be the normalizer of a maximal torus in
\(\operatorname{SL}(2)_k\).

Similarly, even if the image \(H''\) of \(H\) in the finite group
\(G''=G/G^0\) is all of \(G''\), the union of the conjugates of \(H\)
in \(G\) need not be dense. For example, take \(G\) to be the
semidirect product of \(\mathbb Z/2\mathbb Z\) with
\(G^0=\operatorname{SL}(2)_k\), on which it acts by ``symmetry,'' and
take \(H\) to be the semidirect product
\(\mathbb Z/2\mathbb Z\cdot T\), where \(T\) is a maximal torus of
\(G^0\).

On the other hand, if one does not assume in advance that \(H^0\)
contains a Cartan subgroup of \(G^0\), but assumes that the union of
the conjugates of \(H\) in \(G\) is dense, then \(H^0\) necessarily
contains a Cartan subgroup of \(G^0\). To see this, one may plainly
assume that \(G\) is connected and repeat the proof of
\SGARefLink{sga3:xiii:theorem:2:1:item:vi}{2.1(vi)} \(\Rightarrow\) \SGARefLink{sga3:xiii:theorem:2:1:item:i}{(i)}, which does not require \(H\) to be
connected.

% END inputs/XIII/source/tex/components/02_expose_XIII_section1_remainder_and_section2.tex

% BEGIN inputs/XIII/source/tex/components/03_expose_XIII_sections3_4.tex
% Authority: Expo13.pdf, local pp. 14--22 / combined pp. 818--826
% (printed pp. 180--188), from the heading of section 3 through the
% end of Proposition 4.8, stopping immediately before the heading of section 5.
% Authority SHA-256:
% 24735CAFA57291A71C603712C0588358977856D5971503E6E73F3F2E40A70798

\makeatletter
\@ifundefined{remarks}{%
  \theoremstyle{remark}
  \newtheorem{remarks}[theorem]{Remarks}%
}{}
\makeatother

\section{The case of an arbitrary base prescheme}
\SGARefTarget{sga3:xiii:section:3}

Suppose first that the base is a field \(k\), not necessarily
algebraically closed. Since conditions \SGARefLink{sga3:xiii:theorem:2:1}{2.1} \SGARefLink{sga3:xiii:theorem:2:1:item:i-bis}{(i bis)}, \SGARefLink{sga3:xiii:theorem:2:1:item:iv-bis}{(iv bis)}, \SGARefLink{sga3:xiii:theorem:2:1:item:v}{(v)},
and \SGARefLink{sga3:xiii:theorem:2:1:item:vi}{(vi)} are invariant under extension of the base field, passage to
the algebraic closure \(\overline{k}\) of \(k\) shows that they are
equivalent to one another, and equivalent to the assertion that
\(H_{\overline{k}}\) contains a Cartan subgroup of
\(G_{\overline{k}}\). When \SGARefLink{sga3:xiii:theorem:2:1:item:i}{this condition} is satisfied, it is still
true (with the notation of \SGARefLink{sga3:xiii:corollary:2:3}{2.3}) that \(k(X)\) is a finite separable
extension of \(k(G)\), of degree \(d\) independent of every extension
of the base. If \(U\) is the largest open subset of \(G\) such that
\(\psi\) induces a finite étale morphism
\[
  \psi^{-1}(U)\longrightarrow U,
\]
then the formation of \(U\) commutes with extension of the base field.
If \(\psi\) is birational, then \(U\) is also the largest open subset
of \(G\) such that \(\psi\) induces an isomorphism
\[
  \psi^{-1}(U)\xrightarrow{\sim}U;
\]
and if \(g\in U(k)\), there is one and only one subgroup \(H'\) of
\(G\), conjugate to \(H\) over \(\overline{k}\), such that
\(H'\ni g\).

A point \(g\in G(k)\) is called \emph{regular} if it is regular as an
element of
\[
  G(\overline{k})=G_{\overline{k}}(\overline{k}).
\]
More generally, the construction of \SGARefLink{sga3:xiii:corollary:2:7}{2.7} gives an open subset of \(G\),
whose formation commutes with every extension of the base field,
called the \emph{open subset of regular points} of \(G\). It is also
characterized by the following property: for every algebraically
closed extension \(K\) of \(k\) and every \(g\in G(K)\), the point
\(g\) is regular in \(G_K\) if and only if \(g\in U(K)\). If
\(g\in U(k)\)\footnote{Editor's note: \(G(k)\) has been corrected to
\(U(k)\).}, then \(g\) is contained in one and only one Cartan
subgroup of \(G\), as we shall prove below under more general
hypotheses.

Let \(G\) be a smooth, separated group prescheme of finite type, with
connected fibers over the prescheme \(S\), and consider the functor
\SGARefTarget{sga3:xiii:section:3:display:cartan-functor}
\[
  \mathcal C:(\mathrm{Sch})^\circ_{/S}\longrightarrow(\mathrm{Ens})
\]
defined by
\[
  \mathcal C(S')
  =
  \left\{
    \begin{array}{c}
      \text{Cartan subgroups (XII~3.1) of \(G_{S'}\)}
    \end{array}
  \right\}.
\]
Suppose that \SGARefLink{sga3:xiii:section:3:display:cartan-functor}{this functor} is representable by a prescheme smooth over
\(S\). In Exposé XV we give a necessary and sufficient condition for
this to hold, but we already know that the hypothesis is satisfied if
\(G\) is affine over \(S\) with locally constant reductive rank
(\SGARefLink{sga3:xii:3.3}{XII~3.3}), or, more generally, if \(G\) locally admits a maximal torus
for the fpqc topology (\SGARefLink{sga3:xii:7.1}{XII~7.1} a)); this is the case, for example, if
\(S\) is the spectrum of a field.

Let \(X\) be the Cartan subgroup of the \(\mathcal C\)-group
prescheme \(G_{\mathcal C}\), the ``universal Cartan subgroup'' of
\(G\). Thus, as a prescheme over \(S\), \(X\) represents the functor
\[
  X(S')
  =
  \left\{
    \begin{array}{c}
      (C,g)\ \bigm|\ C\text{ is a Cartan subgroup of }G_{S'},\\
      g\text{ is a section of }C\text{ over }S'
    \end{array}
  \right\}.
\]
Consider the canonical projection morphism \((C,g)\mapsto g\),
\[
  \psi:X\longrightarrow G.
\]
We then have \SGARefLink{sga3:xiii:theorem:3:1}{the following theorem}.

\begin{theorem}
\SGARefTarget{sga3:xiii:theorem:3:1}
Under the preceding hypotheses on \(G\), and with the preceding
notation, let \(U\) be the set of \(g\in G\) such that \(g\) is a
regular element of its fiber \(G_s\). Then \(U\) is open, and it is
also the largest open subset \(U\) of \(G\) such that \(\psi\)
induces an isomorphism
\[
  \psi^{-1}(U)\xrightarrow{\sim}U.
\]
\end{theorem}

Let us first prove that \(U\) is open. By the representability of
\(\mathcal C\) as a prescheme smooth over \(S\), and since its
structure morphism is plainly surjective, it follows at once that
\(G\) locally admits a Cartan subgroup for the étale topology and
that the nilpotent rank of the fibers of \(G\) is locally constant.
The same holds for the dimension of the fibers of \(G\). After
localizing on \(S\), we may therefore suppose that these two integers
are constant, say \(\rho\) and \(n\). Consider the Killing polynomial
\[
  P_G(t)=t^n+c_1t^{n-1}+\cdots+c_n\in A[t],
  \qquad A=\Gamma(G,\mathcal O_G).
\]
The restriction of this polynomial to the fibers \(G_s\) of \(G\),
and in particular to the fibers at the maximal points of \(S\), is
divisible by \((t-1)^\rho\). This says that certain linear
combinations of the \(c_i\), with integer coefficients, vanish on the
fibers \(G_s\). When \(S\) is reduced, as we may assume in order to
prove that \(U\) is open, it follows that these combinations
themselves vanish. Hence the Killing polynomial itself is divisible
by \((t-1)^\rho\), say\footnote{The source prints the undefined
exponent \(n-r-1\). It is corrected to \(n-\rho-1\), as required by
the degree \(n-\rho\) of the second factor.}
\[
  P_G(t)
  =
  (t-1)^\rho
  \bigl(t^{n-\rho}+b_1t^{n-\rho-1}+\cdots+b_{n-\rho}\bigr).
\]
Let \(b\) be the sum of the coefficients
\(b_0=1,b_1,\ldots,b_{n-\rho}\) of the second factor. By \SGARefLink{sga3:xiii:corollary:2:7}{2.7}, applied
to the fibers of \(G\), we have
\[
  U=G_b,
\]
which proves that \(U\) is open.

To prove that \(\psi^{-1}(U)\to U\) is an isomorphism, SGA~1,
I~5.7 reduces us to checking this fiber by fiber. We are thereby
reduced to the case of a base field, which may be assumed
algebraically closed. There is then a Cartan subgroup \(C\) of \(G\).
If \(N\) is its normalizer, the conjugacy theorem \SGARefLink{sga3:xii:7.1}{XII~7.1} a) and b)
identifies \(\mathcal C\) with \(G/N\), and the morphism
\(\psi:X\to G\) considered here is precisely the one defined in
\SGARefLink{sga3:xiii:section:2}{Section~2}. The conclusion follows from \SGARefLink{sga3:xiii:theorem:2:6:item:b}{2.6 b)}. The same argument also
shows that \(U\) is the largest open subset of \(G\) for which
\(\psi\) induces an isomorphism \(\psi^{-1}(U)\to U\).

\begin{corollary}
\SGARefTarget{sga3:xiii:corollary:3:2}
Under the hypotheses of \SGARefLink{sga3:xiii:theorem:3:1}{3.1}, let \(g\) be a regular section of \(G\),
that is, suppose that \(g(s)\) is a regular point of \(G_s\) for every
\(s\in S\). Then there is one and only one Cartan subgroup \(C\) of
\(G\) such that \(g\) is a section of \(C\).
\end{corollary}

Indeed, the hypothesis on \(g\) says that \(g\) is a section of
\(U\), while the conclusion says that it admits a unique lift to a
section of \(X\). This is simply another way of saying that
\(\psi^{-1}(U)\to U\) is an isomorphism.

Observe now that the open subset \(\psi^{-1}(U)\) of the Cartan
subgroup \(X\) of \(G_{\mathcal C}\) is precisely the open subset of
\(X\) consisting of the points regular in \(G_{\mathcal C}\), by
which we mean regular in their fibers. We thus obtain a natural
``fibration'' of the dense open subset \(U\) of regular points of
\(G\) over the prescheme \(\mathcal C\); its fibers are dense open
subsets of Cartan subgroups of the fibers of \(G_{\mathcal C}\),
namely their open subsets of points regular in \(G\). For example, we
obtain \SGARefLink{sga3:xiii:corollary:3:3}{the following result}, which will be considerably sharpened in
the next Exposé.

\begin{corollary}
\SGARefTarget{sga3:xiii:corollary:3:3}
Let \(G\) be a smooth connected algebraic group over the field \(k\),
and let \(\mathcal T\) be the scheme of maximal tori of \(G\), which
is isomorphic to the scheme of Cartan subgroups of \(G\). Then the
function field \(k(G)\) of \(G\) is isomorphic to the function field
of a smooth connected affine nilpotent algebraic group \(C\) over the
function field \(k(\mathcal T)\) of \(\mathcal T\), namely
\[
  C=\text{``the generic Cartan subgroup of \(G\)''}.
\]
If \(G\) is affine and has unipotent rank zero, that is, if the Cartan
subgroups of \(G_{\overline k}\)\footnote{The source prints the
undefined base field \(\overline K\); it is corrected to
\(\overline k\).} are tori, then \(k(G)\) is a
unirational extension of \(k(\mathcal T)\).
\end{corollary}

Of course, by the generic Cartan subgroup of \(G\) we mean, by an
abuse of language, the Cartan subgroup of \(G_{k(\mathcal T)}\) that
is the generic fiber of \(X\) over \(\mathcal T\). It remains only to
prove the last assertion of \SGARefLink{sga3:xiii:corollary:3:3}{3.3}, which is contained in \SGARefLink{sga3:xiii:lemma:3:4}{the following
well-known result}, due to Chevalley.

\begin{lemma}
\SGARefTarget{sga3:xiii:lemma:3:4}
Let \(k\) be a field, \(T\) a torus over \(k\), and \(k(T)\) the
field of rational functions on \(T\). Then \(k(T)\) is a
unirational extension of \(k\), that is, it is contained in a purely
transcendental extension of \(k\).
\end{lemma}

Indeed, let \(k'\) be a finite separable extension of \(k\) that
splits\footnote{Editor's note: ``splits'' (?), or: ``over which \(T\)
is split''.} \(T\) (\SGARefLink{sga3:x:proposition:1:4}{X~1.4}). Then \(T\otimes_k k'\) is a rational
variety: it contains a dense open subset isomorphic to a dense open
subset of affine \(n\)-space \(\mathbb A^n_{k'}\). Hence
\[
  T'
  =
  \prod_{\operatorname{Spec}(k')/\operatorname{Spec}(k)}
  T_{k'}/\operatorname{Spec}(k')
\]
is a rational variety, since it has a dense open subset isomorphic to
a dense open subset of
\[
  \prod_{\operatorname{Spec}(k')/\operatorname{Spec}(k)}
  \mathbb A^n_{k'},
\]
which is isomorphic to affine \(mn\)-space over \(k\), where
\(m=[k':k]\). Consider the norm homomorphism from \(T'\) to \(T\),
which is defined whenever \(T\) is a commutative group scheme over
\(k\). The composite
\[
  T\longrightarrow T'\longrightarrow T
\]
is the \(m\)-th power map on \(T\), and is therefore dominant.
Consequently \(T'\to T\) is dominant, proving that \(T\) is
unirational.

Return to the hypotheses of \SGARefLink{sga3:xiii:theorem:3:1}{3.1}, but suppose in addition that \(G\)
locally admits a maximal torus for the fpqc topology (\SGARefLink{sga3:xii:7.1}{XII~7.1}). Let
\(Y\) be the maximal torus of the Cartan subgroup \(X\) of
\(G_{\mathcal C}\). The morphism \(\psi:X\to G\) then induces a
morphism \(Y\to G\) whose set-theoretic image consists of the
semisimple elements of the fibers of \(G\) (XII~8). Finally, \SGARefLink{sga3:xiii:theorem:3:1}{3.1}
shows that the restriction of \(\psi\) to the open subset
\(Y^{\mathrm{reg}}\) of regular points of \(Y\) induces a closed
immersion
\[
  Y^{\mathrm{reg}}\longrightarrow U=G^{\mathrm{reg}}.
\]
Making explicit the functorial meaning of
\(Z=Y^{\mathrm{reg}}\), we obtain \SGARefLink{sga3:xiii:corollary:3:5}{the following}.

\begin{corollary}
\SGARefTarget{sga3:xiii:corollary:3:5}
Let \(G\) be a smooth, separated \(S\)-group prescheme of finite
type, with connected fibers over the prescheme \(S\), and suppose
that \(G\) locally admits a maximal torus for the fpqc topology. Let
\[
  Z:(\mathrm{Sch})^\circ_{/S}\longrightarrow(\mathrm{Ens})
\]
be the functor defined by
\[
  Z(S')
  =
  \left\{
    \begin{array}{c}
      \text{regular sections of \(G_{S'}\) over \(S'\) that are}\\
      \text{contained in a maximal torus of \(G_{S'}\)}
    \end{array}
  \right\}.
\]
Then \(Z\) is representable by a closed subprescheme, smooth over
\(S\), of the open subset \(U=G^{\mathrm{reg}}\) of \(G\) introduced
in \SGARefLink{sga3:xiii:theorem:3:1}{3.1}.
\end{corollary}

To conclude, let us record \SGARefLink{sga3:xiii:corollary:3:6}{the following result}, which sharpens the
density theorem \SGARefLink{sga3:xiii:theorem:2:1}{2.1} \SGARefLink{sga3:xiii:theorem:2:1:item:i}{(i)} \(\Rightarrow\) \SGARefLink{sga3:xiii:theorem:2:1:item:vi}{(vi)}.

\begin{corollary}
\SGARefTarget{sga3:xiii:corollary:3:6}
Under the hypotheses of \SGARefLink{sga3:xiii:corollary:3:5}{3.5}, let \(C\) be a Cartan subgroup of \(G\),
and consider the morphism
\[
  \varphi:Z\times C\longrightarrow G
\]
defined by
\[
  \varphi(g,h)=\operatorname{ad}(g)h=ghg^{-1}.
\]
Then \(\varphi\) is dominant.
\end{corollary}

It plainly suffices to prove the assertion fiber by fiber, which
reduces us to the case where \(S\) is the spectrum of an
algebraically closed field \(k\). Let \(T\) be the maximal torus of
\(C\), let \(t_0\in T(k)\) be regular in \(G\), and let
\(c_0\in C(k)\) be regular in \(G\). Consider
\(\varphi^{-1}(\varphi(t_0,c_0))\). Its \(k\)-rational points are the
pairs \((t,c)\), with \(t\in Z(k)\) and \(c\in C(k)\), such that
\[
  \operatorname{ad}(t)c=\operatorname{ad}(t_0)c_0,
  \qquad\text{that is,}\qquad
  c=\operatorname{ad}(t^{-1}t_0)c_0.
\]
They are therefore in bijection with the \(t\in Z(k)\) such that
\(\operatorname{ad}(t^{-1}t_0)c_0\in C\), or equivalently, since
\(c_0\) is regular, such that \(t^{-1}t_0\in N\), where \(N\) is the
normalizer of \(C\); in other words, \(t\in N\). Restricting to those
\(t\in Z(k)\) that belong to \(C(k)\) gives an open and closed part
of this fiber. We have therefore found a connected component of
\(\varphi^{-1}(\varphi(t_0,c_0))\) isomorphic to \(T\). The preceding
set-theoretic argument indeed gives an isomorphism of schemes: one
simply replaces \(k\)-valued points by points with values in an
arbitrary \(k\)-prescheme.

It follows that the generic fiber of \(\varphi\) has dimension at
most \(\dim T\), and therefore
\[
  \dim\operatorname{Im}(\varphi)
  \geq
  \dim(Z\times C)-\dim T
  =
  \dim Z+\dim C-\dim T.
\]
But
\[
  \dim Z
  =
  \dim Y
  =
  \dim\mathcal C+\dim T
  =
  \dim G-\dim C+\dim T.
\]
Thus \(\dim\operatorname{Im}(\varphi)\geq\dim G\), and \(\varphi\)
is dominant.
\par\hfill\textsc{q.e.d.}

\begin{remarks}
\SGARefTarget{sga3:xiii:remark:3:7}
The argument shows in addition that the connected component at
\((t_0,c_0)\) of the fiber
\(\varphi^{-1}(\varphi(t_0,c_0))\) is isomorphic to \(T\). In
particular, it is smooth over \(k\) and has the same dimension as the
generic fiber. It follows that \(\varphi\) is in fact smooth at
\((t_0,c_0)\), something that one should also be able to verify by
calculating the tangent map. Consequently, under the hypotheses of
\SGARefLink{sga3:xiii:corollary:3:6}{3.6}, the induced morphism
\[
  Z\times_S C^{\mathrm{reg}}\longrightarrow G^{\mathrm{reg}},
  \qquad C^{\mathrm{reg}}=C\cap G^{\mathrm{reg}},
\]
is smooth. Likewise, the analogous morphism
\[
  Z\times T^{\mathrm{reg}}\longrightarrow Z,
\]
where \(T\) is a maximal torus of \(G\), is smooth. More generally,
for every smooth connected normal algebraic subgroup \(H\) of \(C\)
containing an element \(c_0\in G(k)\) that is regular in \(G\), the
image of \(Z\times H\to G\) is dense in the image of
\(G\times H\to G\).
\end{remarks}

\section{Lie algebras over a field: rank, regular elements, Cartan subalgebras}
\SGARefTarget{sga3:xiii:section:4}

In the remainder of \SGARefLink{sga3:xiii:expose}{this Exposé} we take up the theory developed by
Chevalley in his book \emph{Théorie des Groupes de Lie III}
(Act.\ Sc.\ Ind.\ 1226, Paris 1955); the methods of schemes allow us
to dispense with the hypothesis of characteristic zero. We begin by
recalling in \SGARefLink{sga3:xiii:section:4}{this section} several well-known notions and results.

Let \(\mathfrak g\) be a Lie algebra over a ring \(k\). For every
\(a\in\mathfrak g\), denote by \(\operatorname{ad}(a)\) the
endomorphism
\[
  \operatorname{ad}(a)\cdot x=[a,x]
\]
of \(\mathfrak g\); it is a derivation of the Lie algebra
\(\mathfrak g\). For every derivation \(D\) of \(\mathfrak g\), the
nil-space of \(D\), that is, the union of the kernels of the iterates
of \(D\), is a Lie subalgebra of \(\mathfrak g\), as follows from the
Leibniz formula
\[
  D^m([x,y])
  =
  \sum_{0\leq p\leq m}
  \binom{m}{p}[D^p x,D^{m-p}y].
\]
We set
\[
  \operatorname{Nil}(a,\mathfrak g)
  =
  \text{nil-space of }\operatorname{ad}(a)
  =
  \bigcup_{m\geq0}\operatorname{Ker}\operatorname{ad}(a)^m.
\]
When no confusion can arise, we write simply
\(\operatorname{Nil}(a)\), and call it the \emph{nil-space of \(a\)
in \(\mathfrak g\)}.

\begin{proposition}
\SGARefTarget{sga3:xiii:proposition:4:1}
For every \(a\in\mathfrak g\), its nil-space
\(\operatorname{Nil}(a,\mathfrak g)\) is a Lie subalgebra of
\(\mathfrak g\), equal to its own normalizer.
\end{proposition}

It remains only to prove that it is its own normalizer. In other
words, every element of
\(\mathfrak g/\operatorname{Nil}(a,\mathfrak g)\) annihilated by the
adjoint representation of \(\operatorname{Nil}(a,\mathfrak g)\) on
this quotient must be zero. This is immediate, since every element
of the quotient annihilated by \(\operatorname{ad}(a)\)\footnote{The
source prints \(\operatorname{Ad}(a)\); because
\(a\in\mathfrak g\), the endomorphism here is
\(\operatorname{ad}(a)\).} is zero.

For the remainder of \SGARefLink{sga3:xiii:section:4}{this section}, suppose that \(k\) is a field and
that \(\mathfrak g\) is finite-dimensional over \(k\). We denote by
\(W(\mathfrak g)\) the scheme over \(k\) defined by \(\mathfrak g\);
its points in a \(k\)-algebra \(A\) are the elements of
\(\mathfrak g\otimes_k A\). If \(a\in\mathfrak g\), the
characteristic polynomial of \(\operatorname{ad}(a)\) is also called
the \emph{characteristic polynomial}, or the \emph{Killing
polynomial}, of \(a\) in \(\mathfrak g\):
\[
  P_{\mathfrak g}(a,t)
  =
  t^n+c_1(a)t^{n-1}+\cdots+c_n(a),
\]
where \(n=\operatorname{rank}_k\mathfrak g\) and \(c_i(a)\in k\).
Taking the same polynomial for
\(a\in\mathfrak g\otimes_k A\)\footnote{Editor's note:
\(W\otimes_k A\) has been corrected to
\(\mathfrak g\otimes_k A\).}, where \(A\) is any \(k\)-algebra, we
see that the \(c_i(a)\) arise from uniquely determined sections
\(c_i\) of the structure sheaf of \(W(\mathfrak g)\), that is, from
elements of the symmetric algebra
\[
  A=\operatorname{Sym}_k(\mathfrak g^\vee),
\]
where \(\mathfrak g^\vee\) is the dual of the \(k\)-module
\(\mathfrak g\). When \(k\) is infinite, the \(c_i\) are determined
by the corresponding polynomial functions \(\mathfrak g\to k\), but
this is no longer true when \(k\) is finite.

Let \(r\) be the largest integer such that the Killing polynomial
\[
  P_{\mathfrak g}(t)
  =
  t^n+c_1t^{n-1}+\cdots+c_n
  \in A[t]
\]
is divisible by \(t^r\); thus
\[
  P_{\mathfrak g}(t)
  =
  t^n+c_1t^{n-1}+\cdots+c_{n-r}t^r,
  \qquad c_{n-r}\neq0.
\]
The integer \(r\) is called the \emph{nilpotent rank} of the Lie
algebra \(\mathfrak g\). It is invariant under extension of the base
field.

\begin{proposition}
\SGARefTarget{sga3:xiii:proposition:4:2}
Let \(r\) be the nilpotent rank of \(\mathfrak g\), and let
\(a\in\mathfrak g\). Then
\[
  \operatorname{rank}_k\operatorname{Nil}(a,\mathfrak g)\geq r,
\]
with equality if and only if
\[
  c_{n-r}(a)\neq0.
\]
In this case \(\operatorname{Nil}(a,\mathfrak g)\) is a nilpotent
Lie algebra. We shall see the converse in \SGARefLink{sga3:xiii:corollary:5:7:part:b}{5.7 b)} when
\(\mathfrak g\) is the Lie algebra of an algebraic group \(G\)
smooth over \(k\).
\end{proposition}

The first assertion is immediate, since by definition
\[
  \operatorname{rank}_k\operatorname{Nil}(a,\mathfrak g)
  =
  \text{the multiplicity of the zero root in }P_{\mathfrak g}(a,t).
\]
Suppose now that \(c_{n-r}(a)\neq0\). We prove that
\(\operatorname{Nil}(a)\) is nilpotent, which is equivalent to saying
that \(\operatorname{ad}_{\operatorname{Nil}(a)}(x)\) is a nilpotent
endomorphism for every \(x\in\operatorname{Nil}(a)\). We may suppose
that \(k\) is algebraically closed. Since
\[
  \operatorname{ad}(a)_{\mathfrak g/\operatorname{Nil}(a)}
\]
is injective, there is a nonempty open subset
\(U\subset W(\operatorname{Nil}(a))\) such that, for every
\(x\in U(k)\), the endomorphism
\(\operatorname{ad}(x)_{\mathfrak g/\operatorname{Nil}(a)}\) is
injective and hence
\(\operatorname{Nil}(x)\subset\operatorname{Nil}(a)\). We may further
suppose that \(U\) is contained in the open subset on which
\(c_{n-r}\) does not vanish; that open subset is nonempty because
\(c_{n-r}(a)\neq0\). Then
\(\operatorname{Nil}(x)\) and \(\operatorname{Nil}(a)\) have the same
dimension, so
\[
  \operatorname{Nil}(x)=\operatorname{Nil}(a).
\]
It follows that
\(\operatorname{ad}(x)_{\operatorname{Nil}(a)}\) is nilpotent for
every \(x\in U(k)\). By the principle of extension of algebraic
identities, this remains true for every
\(x\in\operatorname{Nil}(a)\), and hence \(\operatorname{Nil}(a)\) is
nilpotent.

An element \(a\in\mathfrak g\) is called \emph{regular} if
\(c_{n-r}(a)\neq0\), that is, if
\[
  \operatorname{rank}_k\operatorname{Nil}(a,\mathfrak g)=r.
\]
When \(k\) is infinite, this also says that
\(\operatorname{rank}_k\operatorname{Nil}(a,\mathfrak g)\) is as
small as possible as \(a\) varies in \(\mathfrak g\). In every case,
the notion of a regular element of \(\mathfrak g\) is invariant under
extension of the base field. The set of regular points of
\(W(\mathfrak g)\), meaning those that arise from regular points with
values in a suitable extension of \(k\), is open, since it is
\[
  W(\mathfrak g)_{c_{n-r}},
\]
the set of points at which \(c_{n-r}\) is invertible.

\begin{corollary}
\SGARefTarget{sga3:xiii:corollary:4:3}
Let \(a\) be a regular element of \(\mathfrak g\), and let
\(\mathfrak h\) be a Lie subalgebra of \(\mathfrak g\) containing
\(a\). Then \(\mathfrak h\) is nilpotent if and only if
\[
  \mathfrak h\subset\operatorname{Nil}(a,\mathfrak g).
\]
In particular, \(\operatorname{Nil}(a,\mathfrak g)\) is a maximal
nilpotent subalgebra of \(\mathfrak g\).
\end{corollary}

Since \(\operatorname{Nil}(a)\) is nilpotent, the inclusion
\(\mathfrak h\subset\operatorname{Nil}(a)\) implies that
\(\mathfrak h\) is nilpotent. Conversely, if \(\mathfrak h\) is
nilpotent, it is contained in the nil-space of its element \(a\);
that is, \(\mathfrak h\subset\operatorname{Nil}(a)\).

\begin{proposition}
\SGARefTarget{sga3:xiii:proposition:4:4}
Suppose that \(k\) is infinite. Let \(\mathfrak d\) be a Lie
subalgebra of \(\mathfrak g\). Consider the following conditions:
\begin{enumerate}
\item[(i)]\SGARefTarget{sga3:xiii:proposition:4:4:item:i}
\(\mathfrak d\) is maximal nilpotent and contains a regular element
of \(\mathfrak g\).
\item[(i bis)]\SGARefTarget{sga3:xiii:proposition:4:4:item:i-bis}
\(\mathfrak d=\operatorname{Nil}(a,\mathfrak g)\) for some regular
element \(a\) of \(\mathfrak g\).
\item[(ii)]\SGARefTarget{sga3:xiii:proposition:4:4:item:ii}
\(\mathfrak d\) is nilpotent and
\(\mathfrak d=\operatorname{Nil}(a,\mathfrak g)\) for some
\(a\in\mathfrak g\).
\item[(ii bis)]\SGARefTarget{sga3:xiii:proposition:4:4:item:ii-bis}
\(\mathfrak d\) is nilpotent, and there exists \(a\in\mathfrak d\)
such that \(\operatorname{ad}(a)_{\mathfrak g/\mathfrak d}\) is
injective.
\item[(iii)]\SGARefTarget{sga3:xiii:proposition:4:4:item:iii}
\(\mathfrak d\) is nilpotent and equal to its own normalizer.
\end{enumerate}
Then
\[
  \SGARefLink{sga3:xiii:proposition:4:4:item:i}{\text{\rm(i)}}
  \Longleftrightarrow
  \text{\rm\SGARefLink{sga3:xiii:proposition:4:4:item:i-bis}{(i bis)}}
  \Longrightarrow
  \SGARefLink{sga3:xiii:proposition:4:4:item:ii}{\text{\rm(ii)}}
  \Longleftrightarrow
  \text{\rm\SGARefLink{sga3:xiii:proposition:4:4:item:ii-bis}{(ii bis)}}
  \Longleftrightarrow
  \SGARefLink{sga3:xiii:proposition:4:4:item:iii}{\text{\rm(iii)}}.
\]
We shall see in \SGARefLink{sga3:xiii:corollary:5:7:part:a}{5.7 a)} that if \(\mathfrak g\) is the Lie algebra of
a smooth algebraic group, then all the preceding conditions are
equivalent.
\end{proposition}

The equivalence of \SGARefLink{sga3:xiii:proposition:4:4:item:i}{(i)} and \SGARefLink{sga3:xiii:proposition:4:4:item:i-bis}{(i bis)} follows immediately from \SGARefLink{sga3:xiii:corollary:4:3}{4.3}, and
these conditions plainly imply \SGARefLink{sga3:xiii:proposition:4:4:item:ii}{(ii)}. The equivalence of \SGARefLink{sga3:xiii:proposition:4:4:item:ii}{(ii)} and
\SGARefLink{sga3:xiii:proposition:4:4:item:ii-bis}{(ii bis)} is also immediate, as is \SGARefLink{sga3:xiii:proposition:4:4:item:ii-bis}{(ii bis)} \(\Rightarrow\) \SGARefLink{sga3:xiii:proposition:4:4:item:iii}{(iii)}
(cf. \SGARefLink{sga3:xiii:proposition:4:1}{4.1}). It remains to prove \SGARefLink{sga3:xiii:proposition:4:4:item:iii}{(iii)} \(\Rightarrow\) \SGARefLink{sga3:xiii:proposition:4:4:item:ii-bis}{(ii bis)}. This
is the only implication that uses the infinitude of \(k\), and it
follows at once from \SGARefLink{sga3:xiii:lemma:4:5}{the next lemma}.

\begin{lemma}
\SGARefTarget{sga3:xiii:lemma:4:5}
Let \(\mathfrak d\) be a nilpotent Lie algebra over an infinite field
\(k\), acting on a finite-dimensional vector space \(V\). Suppose
that the endomorphism \(u(x)\) is not injective for every
\(x\in\mathfrak d\). Then there is a nonzero element \(v\in V\)
annihilated by \(\mathfrak d\).
\end{lemma}

We may suppose that \(k\) is algebraically closed and that
\(\mathfrak d\) is finite-dimensional. Then \(V\) is known to be a
direct sum of finitely many nonzero stable vector subspaces \(V_i\),
\(1\leq i\leq n\), such that, for every \(i\) and every
\(x\in\mathfrak d\), the restriction \(u(x)|_{V_i}\) has a single
eigenvalue \(\lambda_i(x)\) (cf. Bourbaki, \emph{Groupes et Algèbres
de Lie}, Chap.~I, \(\S\)4, Exercise~22). Let \(c_i(x)\) be the
constant term of the characteristic polynomial of
\(u(x)|_{V_i}\), so that
\[
  \lambda_i(x)=0\quad\Longleftrightarrow\quad c_i(x)=0.
\]
Each \(c_i\) is a polynomial function on \(\mathfrak d\), and the
hypothesis says that \(\mathfrak d\) is the union of their zero sets.
Thus one of the \(c_i\) is identically zero. Replacing \(V\) by
\(V_i\), we are reduced to the case in which all \(u(x)\),
\(x\in\mathfrak d\), are nilpotent. Engel's theorem (Bourbaki,
loc.\ cit., Th.~1) then implies that some nonzero \(v\in V\) is
annihilated by \(\mathfrak d\).
\par\hfill\textsc{q.e.d.}

It is easy to see, still assuming that \(k\) is infinite, that
conditions \SGARefLink{sga3:xiii:proposition:4:4:item:i}{(i)} and \SGARefLink{sga3:xiii:proposition:4:4:item:i-bis}{(i bis)} of \SGARefLink{sga3:xiii:proposition:4:4}{4.4} are invariant under every extension
of the base field. When they hold, \(\mathfrak d\) is called a
\emph{Cartan subalgebra} of \(\mathfrak g\). In general, with \(k\)
not necessarily infinite, \(\mathfrak d\) is called a Cartan
subalgebra of \(\mathfrak g\) if it becomes one after one, and hence
every, extension \(k\to k'\) with \(k'\) infinite. It follows in
particular that \(\mathfrak d\) is nilpotent and equal to its own
normalizer.

\begin{proposition}
\SGARefTarget{sga3:xiii:proposition:4:6}
\SGARefTarget{sga3:xiii:proposition:4:6:item:a}
\textup{a)} Let \(a\in\mathfrak g\). If \(a\) is regular, then it is
contained in one and only one Cartan subalgebra of \(\mathfrak g\).
We shall see the converse in \SGARefLink{sga3:xiii:theorem:6:1:part:d}{6.1 d)}, when \(k\) is algebraically
closed and \(\mathfrak g\) is the Lie algebra of a smooth algebraic
group.

\SGARefTarget{sga3:xiii:proposition:4:6:item:b}
\textup{b)} Let \(\mathfrak d\) be a Cartan subalgebra of
\(\mathfrak g\), and let \(a\in\mathfrak d\). Then \(a\) is regular
in \(\mathfrak g\) if and only if
\(\operatorname{ad}(a)_{\mathfrak g/\mathfrak d}\) is injective.
\end{proposition}

For \SGARefLink{sga3:xiii:proposition:4:6:item:a}{a)}, if \(a\) is regular, then
\(\operatorname{Nil}(a,\mathfrak g)\) is a Cartan subalgebra of
\(\mathfrak g\), since this is true after an infinite extension
\(k'\) of \(k\). It follows at once from \SGARefLink{sga3:xiii:corollary:4:3}{4.3} that every Cartan
subalgebra of \(\mathfrak g\) containing \(a\) is equal to this one.
For \SGARefLink{sga3:xiii:proposition:4:6:item:b}{b)}, observe that the nullity\footnote{Editor's note: that is, the
dimension of the nil-space.} of
\(\operatorname{ad}(a)_{\mathfrak g}\) is the sum of the nullities of
\(\operatorname{ad}(a)_{\mathfrak d}\) and
\(\operatorname{ad}(a)_{\mathfrak g/\mathfrak d}\). Since the first
is \(r\), their sum is \(r\) if and only if
\(\operatorname{ad}(a)_{\mathfrak g/\mathfrak d}\) is injective.
\par\hfill\textsc{q.e.d.}

\begin{corollary}
\SGARefTarget{sga3:xiii:corollary:4:7}
Let \(a\) be a regular element of \(\mathfrak g\), let
\(\mathfrak d\) be a Cartan subalgebra of \(\mathfrak g\) containing
\(a\), let \(A\) be a \(k\)-algebra, and let
\(\mathfrak g_A\) and \(\mathfrak d_A\subset\mathfrak g_A\) be the
\(A\)-Lie algebras obtained from \(\mathfrak g\) and \(\mathfrak d\)
by base change; write \(a_A\) for the image of \(a\) in
\(\mathfrak g_A\). Let \(u\) be an automorphism of
\(\mathfrak g_A\). Then
\[
  u(\mathfrak d_A)=\mathfrak d_A
\]
if and only if
\[
  u(a_A)\in\mathfrak d_A.
\]
\end{corollary}

The condition is plainly necessary; let us prove that it is
sufficient. Suppose it holds. Then
\[
  \mathfrak d'=u^{-1}(\mathfrak d_A)
\]
is a Lie subalgebra containing \(a_A\), and every \(b\in\mathfrak d'\)
is such that \(\operatorname{ad}(b)_{\mathfrak d'}\) is
nilpotent.\footnote{The source prints
\(\mathfrak d'=u(\mathfrak d_A)\) and
\(\operatorname{ad}(b)_{\mathfrak d}\). Under the hypothesis
\(u(a_A)\in\mathfrak d_A\), the inverse image
\(u^{-1}(\mathfrak d_A)\) is the subalgebra that contains \(a_A\),
and the nilpotence assertion concerns its restriction
\(\operatorname{ad}(b)_{\mathfrak d'}\).}
Indeed, \(\mathfrak d'\) is isomorphic to \(\mathfrak d_A\), which
has this property, as follows immediately from the definition of
``nilpotent'' in Bourbaki, \emph{Groupes et Algèbres de Lie},
Chap.~I, \(\S\)4, Def.~1. Taking \(b=a_A\), we see that
\(\operatorname{Nil}(b,\mathfrak g_A)\) contains \(\mathfrak d'\).
On the other hand,
\[
  \operatorname{Nil}(b,\mathfrak g_A)=\mathfrak d_A.
\]
Moreover, \(\mathfrak d'\) is locally a direct summand of
\(\mathfrak g_A\), since \(\mathfrak d\) is, and hence is locally a
direct summand of \(\mathfrak d_A\). It is a projective module of the
same rank \(r\) as \(\mathfrak d_A\), and therefore equals
\(\mathfrak d_A\).
\par\hfill\textsc{q.e.d.}

\begin{proposition}
\SGARefTarget{sga3:xiii:proposition:4:8}
Let \(\mathfrak h\) be a Lie subalgebra of \(\mathfrak g\).

\SGARefTarget{sga3:xiii:proposition:4:8:item:a}
\textup{a)} If \(k\) is infinite, the following conditions are
equivalent:
\begin{enumerate}
\item[(i)]\SGARefTarget{sga3:xiii:proposition:4:8:item:a:i}
\(\mathfrak h\) contains a Cartan subalgebra \(\mathfrak d\) of
\(\mathfrak g\).
\item[(ii)]\SGARefTarget{sga3:xiii:proposition:4:8:item:a:ii}
\(\mathfrak h\) contains a regular element \(a\) of \(\mathfrak g\),
and an element \(b\) such that
\(\operatorname{ad}(b)_{\mathfrak g/\mathfrak h}\) is injective.
\item[(iii)]\SGARefTarget{sga3:xiii:proposition:4:8:item:a:iii}
\(\mathfrak h\) has the same nilpotent rank as \(\mathfrak g\), and
contains a regular element of \(\mathfrak g\).
\end{enumerate}
These conditions are invariant under extension of the base field
\(k\).

\SGARefTarget{sga3:xiii:proposition:4:8:item:b}
\textup{b)} Suppose that these conditions hold after a suitable
infinite extension \(k'\) of \(k\). For \(a\in\mathfrak h\), the
element \(a\) is regular in \(\mathfrak g\) if and only if it is
regular in \(\mathfrak h\) and
\(\operatorname{ad}(a)_{\mathfrak g/\mathfrak h}\) is injective;
equivalently, if and only if
\[
  \operatorname{Nil}(a,\mathfrak h)
  =
  \operatorname{Nil}(a,\mathfrak g)
\]
and this common Lie algebra is a Cartan subalgebra of
\(\mathfrak h\).

\SGARefTarget{sga3:xiii:proposition:4:8:item:c}
\textup{c)} Under the hypotheses of \SGARefLink{sga3:xiii:proposition:4:8:item:b}{b)}, let \(\mathfrak d\) be a Lie
subalgebra of \(\mathfrak h\). In order that \(\mathfrak d\) be a
Cartan subalgebra of \(\mathfrak h\), it is sufficient that it be a
Cartan subalgebra of \(\mathfrak g\). We shall see in \SGARefLink{sga3:xiii:corollary:5:8}{5.8} that \SGARefLink{sga3:xiii:proposition:4:8:item:c}{this
condition} is also necessary when \(\mathfrak g\) is the Lie algebra
of a smooth algebraic group \(G\), and \(\mathfrak h\) is the Lie
algebra of a smooth algebraic subgroup \(H\) of \(G\).
\end{proposition}

Conditions \SGARefLink{sga3:xiii:proposition:4:8:item:a:ii}{(ii)} and \SGARefLink{sga3:xiii:proposition:4:8:item:a:iii}{(iii)} of \SGARefLink{sga3:xiii:proposition:4:8:item:a}{a)} are visibly invariant under extension
of the base field \(k\), assumed infinite. In \SGARefLink{sga3:xiii:proposition:4:8:item:b}{b)} and \SGARefLink{sga3:xiii:proposition:4:8:item:c}{c)}, we may
suppose that \(k\) is infinite, and shall do so. If \(\mathfrak h\)
contains the Cartan subalgebra
\[
  \mathfrak d=\operatorname{Nil}(a,\mathfrak g),
\]
then \(a\) is a regular element of \(\mathfrak g\) such that
\(\operatorname{ad}(a)_{\mathfrak g/\mathfrak h}\) is injective.
Thus \SGARefLink{sga3:xiii:proposition:4:8:item:a:i}{(i)} \(\Rightarrow\) \SGARefLink{sga3:xiii:proposition:4:8:item:a:ii}{(ii)}. Conversely, if \SGARefLink{sga3:xiii:proposition:4:8:item:a:ii}{(ii)} holds, then a
``sufficiently general'' element \(a\) of \(\mathfrak h\)
simultaneously satisfies both conditions in \SGARefLink{sga3:xiii:proposition:4:8:item:a:ii}{(ii)}. Hence
\(\operatorname{Nil}(a,\mathfrak g)\) is a Cartan subalgebra of
\(\mathfrak g\) contained in \(\mathfrak h\), proving \SGARefLink{sga3:xiii:proposition:4:8:item:a:i}{(i)}.
Conditions \SGARefLink{sga3:xiii:proposition:4:8:item:a:i}{(i)} and \SGARefLink{sga3:xiii:proposition:4:8:item:a:ii}{(ii)} are therefore equivalent.

Suppose they hold, and let \(a\) vary in \(\mathfrak h\). Then
\SGARefTarget{sga3:xiii:equation:4:8:x}
\[
\tag{x}
  \operatorname{rank}_k\operatorname{Nil}(a,\mathfrak g)
  =
  \operatorname{rank}_k\operatorname{Nil}(a,\mathfrak h)
  +
  \operatorname{rank}_k
  \operatorname{Nil}\bigl(
    \operatorname{ad}(a)_{\mathfrak g/\mathfrak h}
  \bigr).
\]
The two terms on the right are respectively at least \(r'\), the
nilpotent rank of \(\mathfrak h\), and at least \(0\); moreover, both
equalities are attained\footnote{Editor's note: ``the equality'' has
been replaced by ``the equalities''.} for a sufficiently general
element of \(\mathfrak h\). We also have
\[
  \operatorname{rank}_k\operatorname{Nil}(a,\mathfrak g)\geq r,
\]
where \(r\) is the nilpotent rank of \(\mathfrak g\). Equality is
attained for a sufficiently general element of \(\mathfrak h\), and
holds if and only if \(a\) is regular in \(\mathfrak g\).

It follows that \(r=r'\), and that \(a\) is regular if and only if
the two terms on the right-hand side of \(\SGARefLink{sga3:xiii:equation:4:8:x}{\mathrm{(x)}}\) are
respectively \(r'\) and \(0\); that is, if and only if \(a\) is
regular in \(\mathfrak h\) and
\(\operatorname{ad}(a)_{\mathfrak g/\mathfrak h}\) is injective.
This proves \SGARefLink{sga3:xiii:proposition:4:8:item:b}{b)}, and \SGARefLink{sga3:xiii:proposition:4:8:item:c}{c)} follows immediately by taking an element
\(a\in\mathfrak d\) regular in \(\mathfrak g\), so that
\(\operatorname{Nil}(a,\mathfrak g)=\mathfrak d\).

\SGARefLink{sga3:xiii:proposition:4:8:item:b}{The preceding result} also proves \SGARefLink{sga3:xiii:proposition:4:8:item:a:i}{(i)} \(\Rightarrow\) \SGARefLink{sga3:xiii:proposition:4:8:item:a:iii}{(iii)}. Finally,
\SGARefLink{sga3:xiii:proposition:4:8:item:a:iii}{(iii)} \(\Rightarrow\) \SGARefLink{sga3:xiii:proposition:4:8:item:a:i}{(i)}: under \SGARefLink{sga3:xiii:proposition:4:8:item:a:iii}{(iii)}, a sufficiently general
\(a\in\mathfrak h\) is regular both in \(\mathfrak h\) and in
\(\mathfrak g\). Thus
\[
  \operatorname{Nil}(a,\mathfrak h)
  \subset
  \operatorname{Nil}(a,\mathfrak g)
\]
are respectively Cartan subalgebras of \(\mathfrak h\) and
\(\mathfrak g\). Since they have the same rank over \(k\), they are
equal, proving \SGARefLink{sga3:xiii:proposition:4:8:item:a:i}{(i)}. This completes the proof of \SGARefLink{sga3:xiii:proposition:4:8}{4.8}.

% END inputs/XIII/source/tex/components/03_expose_XIII_sections3_4.tex

% BEGIN inputs/XIII/source/tex/components/04_expose_XIII_sections5_6.tex
% Authority: Expo13.pdf, local pp. 22--29 / combined pp. 826--833.
% SHA-256: 24735CAFA57291A71C603712C0588358977856D5971503E6E73F3F2E40A70798.
% The blank local p. 30 / combined p. 834 is excluded; Exposé XIV begins
% on combined p. 835.

\section{The case of the Lie algebra of a smooth algebraic group:
the density theorem}
\SGARefTarget{sga3:xiii:section:5}

Let \(G\) be a smooth algebraic group over the field \(k\), and let
\(\mathfrak g\) be its Lie algebra. Let \(\mathfrak h\) be a Lie
subalgebra of \(\mathfrak g\). Let \(G\) act on \(W(\mathfrak g)\)
through the adjoint representation, and consider the subscheme
\(W(\mathfrak h)\). The construction of \SGARefLink{sga3:xiii:section:1}{Section~1} leads us to introduce
\[
  N=\operatorname{Norm}_G(\mathfrak h)
   =\operatorname{Norm}_G\bigl(W(\mathfrak h)\bigr),
\]
which is an algebraic subgroup of \(G\) (not necessarily smooth), and
\[
  \mathfrak n=\operatorname{Lie}(N),
\]
together with the scheme
\[
  X=G\times^N W(\mathfrak h),
\]
the quotient of \(\widetilde X=G\times W(\mathfrak h)\) by the right
action of \(N\) given by
\[
  (g,x)\cdot n=\bigl(gn,\operatorname{Ad}(n^{-1})x\bigr).
\]
We consider the canonical morphisms
\SGARefTarget{sga3:xiii:section:5:diagram:canonical-morphisms}
\[
\begin{tikzcd}[column sep=large,row sep=large]
G\times W(\mathfrak h)
  \arrow[d,"q"']
  \arrow[dr,"\varphi"] & \\
X \arrow[r,"\psi"] & W(\mathfrak g).
\end{tikzcd}
\]

\begin{theorem}
\SGARefTarget{sga3:xiii:theorem:5:1}
With the preceding notation, suppose that \(k\) is infinite. Consider
the following conditions:
\begin{enumerate}
  \item[(i)]\SGARefTarget{sga3:xiii:theorem:5:1:condition:i}
  \(\mathfrak h\) contains a Cartan subalgebra \(\mathfrak d\) of
  \(\mathfrak g\).

  \item[(ii)]\SGARefTarget{sga3:xiii:theorem:5:1:condition:ii}
  There exists \(a\in\mathfrak h\) such that
  \(\operatorname{ad}(a)_{\mathfrak g/\mathfrak h}\) is injective.

  \item[(iii)]\SGARefTarget{sga3:xiii:theorem:5:1:condition:iii}
  The morphism
  \(\varphi:G\times W(\mathfrak h)\to W(\mathfrak g)\) is generically
  smooth.

  \item[(iv)]\SGARefTarget{sga3:xiii:theorem:5:1:condition:iv}
  \SGARefLink{sga3:xiii:theorem:5:1:condition:iii}{The preceding morphism} \(\varphi\) (or, equivalently,
  \(\psi:X\to W(\mathfrak g)\)) is dominant, and \(\mathfrak h\) has
  the same nilpotent rank as \(\mathfrak g\).

  \item[(v)]\SGARefTarget{sga3:xiii:theorem:5:1:condition:v}
  The morphism \(\psi:X\to W(\mathfrak g)\) is generically smooth and
  \(\mathfrak h=\mathfrak n\).

  \item[(vi)]\SGARefTarget{sga3:xiii:theorem:5:1:condition:vi}
  The morphism \(\psi:X\to W(\mathfrak g)\) is dominant, and
  \(\mathfrak h=\mathfrak n\).

  \item[(vii)]\SGARefTarget{sga3:xiii:theorem:5:1:condition:vii}
  The morphism \(\psi:X\to W(\mathfrak g)\) is dominant.

  \item[(viii)]\SGARefTarget{sga3:xiii:theorem:5:1:condition:viii}
  The morphism \(\psi:X\to W(\mathfrak g)\) is dominant, and \(N\) is
  smooth.

  \item[(ix)]\SGARefTarget{sga3:xiii:theorem:5:1:condition:ix}
  The morphism \(\psi:X\to W(\mathfrak g)\) is dominant,
  \(\mathfrak h=\mathfrak n\), and \(\mathfrak h\) is the Lie algebra
  of a smooth algebraic subgroup \(H\) of \(G\).

  \item[(x)]\SGARefTarget{sga3:xiii:theorem:5:1:condition:x}
  The morphism \(\psi:X\to W(\mathfrak g)\) is generically
  quasi-finite, and \(\mathfrak n=\mathfrak h\).

  \item[(xi)]\SGARefTarget{sga3:xiii:theorem:5:1:condition:xi}
  The morphism \(\psi:X\to W(\mathfrak g)\) is generically étale, and
  \(\mathfrak n=\mathfrak h\).

  \item[(xii)]\SGARefTarget{sga3:xiii:theorem:5:1:condition:xii}
  There exists a smooth algebraic subgroup \(H\) of \(G\), with Lie
  algebra \(\mathfrak h\), such that \(\mathfrak h\) contains a Cartan
  subalgebra \(\mathfrak d\) of \(\mathfrak g\).

  \item[(xiii)]\SGARefTarget{sga3:xiii:theorem:5:1:condition:xiii}
  There exists \(a\in\mathfrak h\) such that \(N\) and the transporter
  \(M_a\) of \(a\) into \(\mathfrak h\) (cf.\ \SGARefLink{sga3:xiii:section:1}{Section~1}) coincide in a
  neighborhood of \(e\), and \(\mathfrak n=\mathfrak h\).
\end{enumerate}
One then has \SGARefLink{sga3:xiii:theorem:5:1:diagram:implications}{the following diagram} of implications:
\SGARefTarget{sga3:xiii:theorem:5:1:diagram:implications}
\[
\begin{tikzcd}[column sep=1.45em,row sep=3.2em]
\SGARefLink{sga3:xiii:theorem:5:1:condition:xi}{(\mathrm{xi})}
  \arrow[rr,Leftrightarrow] &&
\SGARefLink{sga3:xiii:theorem:5:1:condition:xii}{(\mathrm{xii})}
  \arrow[r,Leftrightarrow] &
\SGARefLink{sga3:xiii:theorem:5:1:condition:xiii}{(\mathrm{xiii})}
  \arrow[r,Rightarrow] &
\SGARefLink{sga3:xiii:theorem:5:1:condition:viii}{(\mathrm{viii})}
  \arrow[r,Leftrightarrow] &
\SGARefLink{sga3:xiii:theorem:5:1:condition:ix}{(\mathrm{ix})}
  \arrow[r,Leftrightarrow] &
\SGARefLink{sga3:xiii:theorem:5:1:condition:x}{(\mathrm{x})} \\
\SGARefLink{sga3:xiii:theorem:5:1:condition:i}{(\mathrm{i})}
  \arrow[r,Leftrightarrow] &
\SGARefLink{sga3:xiii:theorem:5:1:condition:ii}{(\mathrm{ii})}
  \arrow[r,Leftrightarrow] &
\SGARefLink{sga3:xiii:theorem:5:1:condition:iii}{(\mathrm{iii})}
  \arrow[u,Leftarrow]
  \arrow[r,Leftrightarrow] &
\SGARefLink{sga3:xiii:theorem:5:1:condition:iv}{(\mathrm{iv})}
  \arrow[r,Rightarrow] &
\SGARefLink{sga3:xiii:theorem:5:1:condition:v}{(\mathrm{v})}
  \arrow[r,Rightarrow] &
\SGARefLink{sga3:xiii:theorem:5:1:condition:vi}{(\mathrm{vi})}
  \arrow[u,Leftarrow]
  \arrow[r,Rightarrow] &
\SGARefLink{sga3:xiii:theorem:5:1:condition:vii}{(\mathrm{vii})}.
\end{tikzcd}
\]
When \(k\) has characteristic zero, all the conditions under
consideration are equivalent. Finally,
\[
  \SGARefLink{sga3:xiii:theorem:5:1:condition:xi}{(\mathrm{xi})}
  \Longleftrightarrow
  [\SGARefLink{sga3:xiii:theorem:5:1:condition:i}{(\mathrm{i})}\ \text{and}\ \SGARefLink{sga3:xiii:theorem:5:1:condition:viii}{(\mathrm{viii})}]
  \Longleftrightarrow
  [\SGARefLink{sga3:xiii:theorem:5:1:condition:v}{(\mathrm{v})}\ \text{and}\ \SGARefLink{sga3:xiii:theorem:5:1:condition:viii}{(\mathrm{viii})}].
\]
\end{theorem}

Let us first record the trivial implications
\[
  \SGARefLink{sga3:xiii:theorem:5:1:condition:v}{(\mathrm{v})}\Rightarrow\SGARefLink{sga3:xiii:theorem:5:1:condition:vi}{(\mathrm{vi})}\Rightarrow\SGARefLink{sga3:xiii:theorem:5:1:condition:vii}{(\mathrm{vii})},
  \qquad
  \SGARefLink{sga3:xiii:theorem:5:1:condition:ix}{(\mathrm{ix})}\Rightarrow\SGARefLink{sga3:xiii:theorem:5:1:condition:vi}{(\mathrm{vi})},
  \qquad
  \SGARefLink{sga3:xiii:theorem:5:1:condition:xi}{(\mathrm{xi})}
  \Longleftrightarrow
  [\SGARefLink{sga3:xiii:theorem:5:1:condition:x}{(\mathrm{x})}\ \text{and}\ \SGARefLink{sga3:xiii:theorem:5:1:condition:v}{(\mathrm{v})}].
\]

We prove the equivalence of conditions \SGARefLink{sga3:xiii:theorem:5:1:condition:i}{(i)} through \SGARefLink{sga3:xiii:theorem:5:1:condition:iv}{(iv)}, and that they
imply \SGARefLink{sga3:xiii:theorem:5:1:condition:v}{(v)}. The implication \SGARefLink{sga3:xiii:theorem:5:1:condition:i}{(i)}\(\Rightarrow\)\SGARefLink{sga3:xiii:theorem:5:1:condition:ii}{(ii)} is immediate. On the
other hand, when \(k\) is algebraically closed, condition \SGARefLink{sga3:xiii:theorem:5:1:condition:iii}{(iii)} says
that there is a \(k\)-rational point of \(G\times W(\mathfrak h)\) at
which the tangent map of \(\varphi\) is surjective. It is clear that,
after translating by a suitable element of \(G(k)\), this point may be
taken to have the form \((e,a)\), with \(a\in\mathfrak h\). It follows
that when \(k\) is infinite, though not necessarily algebraically
closed, \SGARefLink{sga3:xiii:theorem:5:1:condition:iii}{this condition}---which is plainly sufficient for generic
smoothness---is also necessary.

The tangent map is readily computed. Upon identifying the tangent
space of \(W(\mathfrak h)\) at \(a\) with \(\mathfrak h\), it is the
map
\[
  (\xi,x)\longmapsto[\xi,a]+x
\]
from \(\mathfrak g\times\mathfrak h\) to \(\mathfrak g\). Its
surjectivity is equivalent to the surjectivity of
\(\operatorname{ad}(a)_{\mathfrak g/\mathfrak h}\), or, what amounts
to the same thing, to its injectivity. This proves the equivalence of
\SGARefLink{sga3:xiii:theorem:5:1:condition:ii}{(ii)} and \SGARefLink{sga3:xiii:theorem:5:1:condition:iii}{(iii)}. Moreover, \SGARefLink{sga3:xiii:theorem:5:1:condition:ii}{(ii)} plainly implies
\(\mathfrak h=\mathfrak n\), and \SGARefLink{sga3:xiii:theorem:5:1:condition:iii}{(iii)} implies that \(\psi\) is
generically smooth: if \(\varphi\) is smooth at a point \(u\), then,
since \(q\) is flat, \(\psi\) is smooth at \(q(u)\). Thus \SGARefLink{sga3:xiii:theorem:5:1:condition:ii}{(ii)} and
\SGARefLink{sga3:xiii:theorem:5:1:condition:iii}{(iii)} imply \SGARefLink{sga3:xiii:theorem:5:1:condition:v}{(v)}.

We now show that they imply \SGARefLink{sga3:xiii:theorem:5:1:condition:i}{(i)}. Since \(\psi\) is dominant and the set
of regular points of \(\mathfrak g\) is dense and open, \(\mathfrak h\)
contains regular elements of \(\mathfrak g\). Hence a sufficiently
general element \(b\) of \(\mathfrak h\) is regular in \(\mathfrak g\)
and makes \(\operatorname{ad}(b)_{\mathfrak g/\mathfrak h}\)
injective. Consequently
\(\operatorname{Nil}(b,\mathfrak g)\subset\mathfrak h\), so
\(\mathfrak h\) contains the Cartan subalgebra
\(\mathfrak d=\operatorname{Nil}(b,\mathfrak g)\). Conditions \SGARefLink{sga3:xiii:theorem:5:1:condition:i}{(i)},
\SGARefLink{sga3:xiii:theorem:5:1:condition:ii}{(ii)}, and \SGARefLink{sga3:xiii:theorem:5:1:condition:iii}{(iii)} are therefore equivalent. Finally,
\SGARefLink{sga3:xiii:theorem:5:1:condition:i}{(i)}\(\Longleftrightarrow\)\SGARefLink{sga3:xiii:theorem:5:1:condition:iv}{(iv)}: we have already observed that if
\(\mathfrak h\) contains a Cartan subalgebra, then it has the same
nilpotent rank as \(\mathfrak g\) (\SGARefLink{sga3:xiii:proposition:4:6}{4.6}), so \SGARefLink{sga3:xiii:theorem:5:1:condition:i}{(i)} implies \SGARefLink{sga3:xiii:theorem:5:1:condition:iv}{(iv)}.
Conversely, if \SGARefLink{sga3:xiii:theorem:5:1:condition:iv}{(iv)} holds, then \(\mathfrak h\) contains a regular
element of \(\mathfrak g\), and since it has the same rank as
\(\mathfrak g\), it contains a Cartan subalgebra by \SGARefLink{sga3:xiii:proposition:4:6}{4.6}.

We next prove the equivalence of conditions \SGARefLink{sga3:xiii:theorem:5:1:condition:viii}{(viii)} through \SGARefLink{sga3:xiii:theorem:5:1:condition:x}{(x)}. We
first note \SGARefLink{sga3:xiii:lemma:5:2}{the following facts}.

\begin{lemma}
\SGARefTarget{sga3:xiii:lemma:5:2}
\begin{enumerate}
  \item[a)]\SGARefTarget{sga3:xiii:lemma:5:2:part:a}
  If \(N\) is smooth, then
  \[
    \dim X\leq\dim G,
  \]
  with equality if and only if \(\mathfrak h=\mathfrak n\).

  \item[b)]\SGARefTarget{sga3:xiii:lemma:5:2:part:b}
  If \(\mathfrak h=\mathfrak n\), then
  \[
    \dim X\geq\dim G,
  \]
  with equality if and only if \(N\) is smooth.
\end{enumerate}
\end{lemma}

These assertions follow at once from \SGARefLink{sga3:xiii:lemma:5:2:display:dimension-formula}{the formula}
\SGARefTarget{sga3:xiii:lemma:5:2:display:dimension-formula}
\[
  \dim X=\dim G-\dim N+\dim_k\mathfrak h
\]
and from the fact that
\(\dim N=\dim_k\mathfrak n\) is equivalent to the smoothness of \(N\).

Now \SGARefLink{sga3:xiii:theorem:5:1:condition:viii}{(viii)} implies \SGARefLink{sga3:xiii:theorem:5:1:condition:x}{(x)}: by dominance, \SGARefLink{sga3:xiii:theorem:5:1:condition:viii}{(viii)} gives
\(\dim X\geq\dim G\), so \SGARefLink{sga3:xiii:lemma:5:2:part:a}{Lemma~5.2(a)} gives equality and
\(\mathfrak h=\mathfrak n\), hence \SGARefLink{sga3:xiii:theorem:5:1:condition:x}{(x)}. The same argument in reverse,
using \SGARefLink{sga3:xiii:lemma:5:2:part:b}{Lemma~5.2(b)}, shows that \SGARefLink{sga3:xiii:theorem:5:1:condition:x}{(x)} implies \SGARefLink{sga3:xiii:theorem:5:1:condition:viii}{(viii)}. Moreover, \SGARefLink{sga3:xiii:theorem:5:1:condition:viii}{(viii)}
implies \SGARefLink{sga3:xiii:theorem:5:1:condition:ix}{(ix)}, since it implies \(\mathfrak h=\mathfrak n\), so
\(\mathfrak h\) is the Lie algebra of the smooth algebraic subgroup
\(N\) of \(G\). Conversely, \SGARefLink{sga3:xiii:theorem:5:1:condition:ix}{(ix)} implies \SGARefLink{sga3:xiii:theorem:5:1:condition:viii}{(viii)}: since \(H\)
normalizes its Lie algebra \(\mathfrak h\), it is contained in \(N\);
and since \(H\) is smooth and has the same Lie algebra as \(N\), it
follows that \(N\) is smooth.

Finally, we prove the equivalence of \SGARefLink{sga3:xiii:theorem:5:1:condition:xi}{(xi)}, \SGARefLink{sga3:xiii:theorem:5:1:condition:xii}{(xii)}, and \SGARefLink{sga3:xiii:theorem:5:1:condition:xiii}{(xiii)}, and show
that they imply \SGARefLink{sga3:xiii:theorem:5:1:condition:iii}{(iii)}; this will complete \SGARefLink{sga3:xiii:theorem:5:1:diagram:implications}{the implication diagram}.
We have \SGARefLink{sga3:xiii:theorem:5:1:condition:xi}{(xi)}\(\Longleftrightarrow\)\SGARefLink{sga3:xiii:theorem:5:1:condition:xiii}{(xiii)}. Indeed, if
\(\mathfrak n=\mathfrak h\), then \SGARefLink{sga3:xiii:lemma:5:2:part:b}{Lemma~5.2(b)} gives
\(\dim X\geq\dim G\). Since \(W(\mathfrak g)\) is normal, condition
\SGARefLink{sga3:xiii:theorem:5:1:condition:xi}{(xi)} is therefore equivalent to \(\psi\) being generically
unramified, and this in turn is equivalent to \SGARefLink{sga3:xiii:theorem:5:1:condition:xiii}{(xiii)} by
\SGARefLink{sga3:xiii:lemma:1:1}{1.1}, \SGARefLink{sga3:xiii:theorem:5:1:condition:ii}{(ii)}\(\Longleftrightarrow\)\SGARefLink{sga3:xiii:theorem:5:1:condition:iii}{(iii)}, arguing as above in the proof
of \SGARefLink{sga3:xiii:theorem:5:1:condition:ii}{(ii)}\(\Longleftrightarrow\)\SGARefLink{sga3:xiii:theorem:5:1:condition:iii}{(iii)}.

Since \SGARefLink{sga3:xiii:theorem:5:1:condition:xi}{(xi)}\(\Rightarrow\)\SGARefLink{sga3:xiii:theorem:5:1:condition:x}{(x)}\(\Rightarrow\)\SGARefLink{sga3:xiii:theorem:5:1:condition:viii}{(viii)}, what we have
already proved shows that \SGARefLink{sga3:xiii:theorem:5:1:condition:xi}{(xi)} implies that \(N\) is smooth. Thus
\(q:G\times W(\mathfrak h)\to X\) is smooth, and the composite
\(\varphi=\psi\circ q\) is generically smooth; this is \SGARefLink{sga3:xiii:theorem:5:1:condition:iii}{(iii)}. Since
\SGARefLink{sga3:xiii:theorem:5:1:condition:iii}{(iii)} implies \SGARefLink{sga3:xiii:theorem:5:1:condition:i}{(i)}, it also follows that
\SGARefLink{sga3:xiii:theorem:5:1:condition:xi}{(xi)}\(\Rightarrow\)\SGARefLink{sga3:xiii:theorem:5:1:condition:xii}{(xii)}.\footnote{Editor's note: because \(H=N\)
will do.} Finally, \SGARefLink{sga3:xiii:theorem:5:1:condition:xii}{(xii)} implies \SGARefLink{sga3:xiii:theorem:5:1:condition:xi}{(xi)}. Indeed, \SGARefLink{sga3:xiii:theorem:5:1:condition:xii}{(xii)} plainly implies
\SGARefLink{sga3:xiii:theorem:5:1:condition:i}{(i)}, and we have seen that
\SGARefLink{sga3:xiii:theorem:5:1:condition:i}{(i)}\(\Rightarrow\)\SGARefLink{sga3:xiii:theorem:5:1:condition:iii}{(iii)}\(\Rightarrow\)\SGARefLink{sga3:xiii:theorem:5:1:condition:v}{(v)}; hence \SGARefLink{sga3:xiii:theorem:5:1:condition:xii}{(xii)} implies \SGARefLink{sga3:xiii:theorem:5:1:condition:v}{(v)}.
It also implies \SGARefLink{sga3:xiii:theorem:5:1:condition:ix}{(ix)}, and \SGARefLink{sga3:xiii:theorem:5:1:condition:ix}{(ix)} implies \SGARefLink{sga3:xiii:theorem:5:1:condition:x}{(x)}. Thus \SGARefLink{sga3:xiii:theorem:5:1:condition:xii}{(xii)} implies both
\SGARefLink{sga3:xiii:theorem:5:1:condition:v}{(v)} and \SGARefLink{sga3:xiii:theorem:5:1:condition:x}{(x)}, and therefore \SGARefLink{sga3:xiii:theorem:5:1:condition:xi}{(xi)}, since
\[
  \text{generically étale}
  =
  \text{generically smooth}
  +
  \text{generically quasi-finite}.
\]

If \(k\) has characteristic zero, then \SGARefLink{sga3:xiii:theorem:5:1:condition:vii}{(vii)} implies \SGARefLink{sga3:xiii:theorem:5:1:condition:viii}{(viii)}, because
Cartier's theorem makes \(N\) automatically smooth
(VI\textsubscript{B}~1.6.1). Moreover, [\SGARefLink{sga3:xiii:theorem:5:1:condition:viii}{(viii)} and \SGARefLink{sga3:xiii:theorem:5:1:condition:x}{(x)}] implies \SGARefLink{sga3:xiii:theorem:5:1:condition:xi}{(xi)},
because in characteristic zero, for a morphism of integral
preschemes,
\[
  \text{generically étale}
  =
  \text{dominant and generically quasi-finite}.
\]
Thus in this case all conditions \SGARefLink{sga3:xiii:theorem:5:1:condition:i}{(i)} through \SGARefLink{sga3:xiii:theorem:5:1:condition:xiii}{(xiii)} are equivalent.

\begin{corollary}
\SGARefTarget{sga3:xiii:corollary:5:3}
Under the equivalent conditions \SGARefLink{sga3:xiii:theorem:5:1:condition:viii}{(viii)} through \SGARefLink{sga3:xiii:theorem:5:1:condition:x}{(x)}, there exists a
unique smooth connected algebraic subgroup \(H\) of \(G\) whose Lie
algebra is \(\mathfrak h\), and
\[
  \operatorname{Norm}_G(H)
  =\operatorname{Norm}_G(\mathfrak h)
  =N,
  \qquad
  H=N^0.
\]
\end{corollary}

Indeed, \(H=N^0\) has the required properties. Conversely, if \(H\)
has these properties, then \(H\subset N\), since \(H\) normalizes its
Lie algebra \(\mathfrak h\). This is an inclusion of smooth groups
with the same Lie algebra, and \(H\) is connected, so \(H=N^0\). To
prove
\(\operatorname{Norm}_G(H)=\operatorname{Norm}_G(\mathfrak h)\), we
may assume that \(k\) is algebraically closed. What we have just proved
then shows immediately that the two groups have the same \(k\)-points.
On the other hand, the inclusions
\[
  H\subset\operatorname{Norm}_G(H)\subset N
\]
show that \(\operatorname{Norm}_G(H)\) and \(N\) have the same Lie
algebra; hence they are equal.

\begin{corollary}
\SGARefTarget{sga3:xiii:corollary:5:4}
Under the equivalent conditions \SGARefLink{sga3:xiii:theorem:5:1:condition:i}{(i)} through \SGARefLink{sga3:xiii:theorem:5:1:condition:iv}{(iv)}, let
\(a\in\mathfrak h\). Then the following conditions are equivalent, and
they hold whenever \(a\) is regular in \(\mathfrak g\):
\begin{enumerate}
  \item[(i)]\SGARefTarget{sga3:xiii:corollary:5:4:condition:i}
  \(\varphi\) is smooth at \((e,a)\).

  \item[(ii)]\SGARefTarget{sga3:xiii:corollary:5:4:condition:ii}
  \(M_a=\operatorname{Transp}_G(a,\mathfrak h)\) is smooth at \(e\),
  and
  \(\dim_e(M_a)=\operatorname{rank}_k\mathfrak h\).

  \item[(iii)]\SGARefTarget{sga3:xiii:corollary:5:4:condition:iii}
  \(\operatorname{ad}(a)_{\mathfrak g/\mathfrak h}\) is injective
  (equivalently, bijective).
\end{enumerate}
Under the equivalent conditions \SGARefLink{sga3:xiii:theorem:5:1:condition:xi}{(xi)} through \SGARefLink{sga3:xiii:theorem:5:1:condition:xiii}{(xiii)}, let \(H\) be the
algebraic subgroup of \(G\) considered in \SGARefLink{sga3:xiii:corollary:5:3}{Corollary~5.3}. The preceding
conditions are also equivalent to:
\begin{enumerate}
  \item[(iv)]\SGARefTarget{sga3:xiii:corollary:5:4:condition:iv}
  \(\psi\) is étale at \((\overline e,a)\).

  \item[(v)]\SGARefTarget{sga3:xiii:corollary:5:4:condition:v}
  If \(M_a^0\) denotes the connected component of \(e\) in the
  transporter \(M_a\) of \(a\) into \(\mathfrak h\), endowed with the
  structure induced by \(M_a\), then
  \[
    H=M_a^0.
  \]
\end{enumerate}
\end{corollary}

Plainly \SGARefLink{sga3:xiii:corollary:5:4:condition:i}{(i)} implies \SGARefLink{sga3:xiii:corollary:5:4:condition:ii}{(ii)}, since \(M_a\) is isomorphic to the fiber
\(\varphi^{-1}(a)\), with \(e\) corresponding to \((e,a)\).
Conversely, \SGARefLink{sga3:xiii:corollary:5:4:condition:ii}{(ii)} implies \SGARefLink{sga3:xiii:corollary:5:4:condition:i}{(i)}: condition \SGARefLink{sga3:xiii:corollary:5:4:condition:ii}{(ii)} says that \(\varphi\) is
``equidimensional'' at \((e,a)\), in the sense that the dimension of
the fiber through this point is the dimension of the generic fiber.
Since \(G\times W(\mathfrak h)\) and \(W(\mathfrak g)\) are regular,
this implies that \(\varphi\) is flat at \((e,a)\), hence smooth there
because its fiber is smooth at that point. The equivalence of \SGARefLink{sga3:xiii:corollary:5:4:condition:i}{(i)} and
\SGARefLink{sga3:xiii:corollary:5:4:condition:iii}{(iii)} was obtained in the proof of \SGARefLink{sga3:xiii:theorem:5:1}{Theorem~5.1} by the elementary
calculation of the tangent map. We also saw in \SGARefLink{sga3:xiii:proposition:4:8:item:b}{4.8(b)} that regularity
of \(a\) in \(\mathfrak g\) implies \SGARefLink{sga3:xiii:corollary:5:4:condition:iii}{(iii)}.

Under conditions \SGARefLink{sga3:xiii:theorem:5:1:condition:xi}{(xi)} through \SGARefLink{sga3:xiii:theorem:5:1:condition:xiii}{(xiii)}, the morphism
\(q:G\times W(\mathfrak h)\to X\) is smooth, since \(N\) is smooth.
Thus \SGARefLink{sga3:xiii:corollary:5:4:condition:i}{(i)} is equivalent to the smoothness of \(\psi\) at
\((\overline e,a)\); and since \(\psi\) is generically étale, this is
equivalent to \SGARefLink{sga3:xiii:corollary:5:4:condition:iv}{(iv)}. Finally, as was noted at the end of the proof of
\SGARefLink{sga3:xiii:lemma:1:1}{1.1}, condition \SGARefLink{sga3:xiii:corollary:5:4:condition:iv}{(iv)} implies that \(N\) is the prescheme induced on
\(M_a\) by an open and closed subset of \(M_a\), which gives \SGARefLink{sga3:xiii:corollary:5:4:condition:v}{(v)}.
Conversely, \SGARefLink{sga3:xiii:corollary:5:4:condition:v}{(v)} trivially implies \SGARefLink{sga3:xiii:corollary:5:4:condition:ii}{(ii)}. Alternatively, \SGARefLink{sga3:xiii:corollary:5:4:condition:v}{(v)} implies
\SGARefLink{sga3:xiii:corollary:5:4:condition:iv}{(iv)} by \SGARefLink{sga3:xiii:lemma:1:1}{1.1}, since \(\psi\) is dominant and \(W(\mathfrak g)\) is
normal, and ``unramified'' is equivalent here to ``étale''. This
completes the proof of \SGARefLink{sga3:xiii:corollary:5:4}{Corollary~5.4}.

\begin{corollary}
\SGARefTarget{sga3:xiii:corollary:5:5}
Let \(G\) be a smooth algebraic group over a field \(k\), and let
\(H\) be a smooth algebraic subgroup such that its Lie algebra
\(\mathfrak h\) contains, at least after a suitable extension of
\(k\), a Cartan subalgebra of the Lie algebra \(\mathfrak g\) of
\(G\). Let \(K\) be a connected algebraic subgroup of \(G\), not
necessarily smooth, with Lie algebra \(\mathfrak k\), and suppose that
\(\mathfrak k\) contains a regular element \(a\) of \(\mathfrak g\),
again at least after a suitable extension of \(k\). Then \(H\)
contains \(K\) if and only if \(\mathfrak h\) contains
\(\mathfrak k\).
\end{corollary}

Indeed, \SGARefLink{sga3:xiii:corollary:5:4}{Corollary~5.4}, \SGARefLink{sga3:xiii:corollary:5:4:condition:iii}{(iii)}\(\Rightarrow\)\SGARefLink{sga3:xiii:corollary:5:4:condition:v}{(v)}, applied to the
connected subgroup supplied by \SGARefLink{sga3:xiii:corollary:5:3}{Corollary~5.3}, gives
\(H^0=M_a^0\).\footnote{The French source prints \(H=M_a^0\) here. In the stated
generality the argument gives \(H^0=M_a^0\), which is all that is
needed.}
This identity is invariant under extension of the base field, so it
remains valid without assuming that the base field is infinite.
Moreover, \(\mathfrak k\subset\mathfrak h\) plainly implies
\(K\subset M_a\). Since \(K\) is connected,
\(K\subset M_a^0\subset H\). \hfill\(\square\)

\begin{corollary}
\SGARefTarget{sga3:xiii:corollary:5:6}
Let \(G\) and \(H\) be as in \SGARefLink{sga3:xiii:corollary:5:5}{Corollary~5.5}, and let \(K\) be an
algebraic subgroup of \(G\). Suppose that \(H\) is connected and \(K\)
is smooth. Then \(K\) contains \(H\) if and only if
\(\mathfrak k\) contains \(\mathfrak h\).
\end{corollary}

Indeed, if \(\mathfrak k\supset\mathfrak h\), then \(K\) satisfies,
with respect to \(H\), the hypothesis imposed on \(H\) in
\SGARefLink{sga3:xiii:corollary:5:5}{Corollary~5.5}. Meanwhile \(H\) plainly satisfies, with respect to
\(K\), the condition imposed on \(K\) there. The conclusion follows
from \SGARefLink{sga3:xiii:corollary:5:5}{Corollary~5.5}.

\begin{corollary}
\SGARefTarget{sga3:xiii:corollary:5:7}
Let \(\mathfrak g\) be the Lie algebra of a smooth algebraic group
\(G\) over a field \(k\). Then:
\begin{enumerate}
  \item[a)]\SGARefTarget{sga3:xiii:corollary:5:7:part:a}
  A Lie subalgebra \(\mathfrak d\) of \(\mathfrak g\) is a Cartan
  subalgebra if and only if it is nilpotent and equal to its own
  normalizer.

  \item[b)]\SGARefTarget{sga3:xiii:corollary:5:7:part:b}
  An element \(a\in\mathfrak g\) is regular if and only if
  \(\operatorname{Nil}(a,\mathfrak g)\) is nilpotent.
\end{enumerate}
\end{corollary}

After extending the base field, we may assume that \(k\) is infinite.
In view of \SGARefLink{sga3:xiii:proposition:4:4}{4.4}, to prove \SGARefLink{sga3:xiii:corollary:5:7:part:a}{(a)} it remains only to show that if
\(\mathfrak d\) is nilpotent and contains an element \(a\) for which
\(\operatorname{ad}(a)_{\mathfrak g/\mathfrak d}\) is injective, then
\(\mathfrak d\) is a Cartan subalgebra. By \SGARefLink{sga3:xiii:theorem:5:1}{Theorem~5.1},
\SGARefLink{sga3:xiii:theorem:5:1:condition:ii}{(ii)}\(\Rightarrow\)\SGARefLink{sga3:xiii:theorem:5:1:condition:i}{(i)}, \(\mathfrak d\) contains a Cartan subalgebra
\[
  \mathfrak d'=\operatorname{Nil}(a,\mathfrak g).
\]
Since \(\mathfrak d\) is nilpotent, \SGARefLink{sga3:xiii:corollary:4:3}{4.3} now gives
\(\mathfrak d=\mathfrak d'\). For \SGARefLink{sga3:xiii:corollary:5:7:part:b}{(b)}, one notes that
\(\operatorname{Nil}(a,\mathfrak g)\) is a Cartan subalgebra of
\(\mathfrak g\) by \SGARefLink{sga3:xiii:corollary:5:7:part:a}{(a)}, and therefore \(a\) is regular.

\begin{corollary}
\SGARefTarget{sga3:xiii:corollary:5:8}
Let \(G\) be a smooth algebraic group over a field \(k\), and let
\(H\) be a smooth algebraic subgroup, with Lie algebras
\(\mathfrak g\) and \(\mathfrak h\), respectively. Suppose that after
a suitable extension of the base field, \(\mathfrak h\) contains a
Cartan subalgebra of \(\mathfrak g\). Let \(\mathfrak d\) be a Lie
subalgebra of \(\mathfrak h\). Then \(\mathfrak d\) is a Cartan
subalgebra of \(\mathfrak h\) if and only if it is a Cartan subalgebra
of \(\mathfrak g\).
\end{corollary}

In view of \SGARefLink{sga3:xiii:proposition:4:8:item:c}{4.8(c)}, it remains to show that if \(\mathfrak d\) is a
Cartan subalgebra of \(\mathfrak h\), then it is one of
\(\mathfrak g\). It is enough to show that \(\mathfrak d\) contains
an element \(a\) regular in \(\mathfrak g\), and we may assume that
\(k\) is algebraically closed. But there is a dense open subset of
\(\mathfrak h\) consisting of points regular in \(\mathfrak g\), so
the assertion follows from \SGARefLink{sga3:xiii:theorem:5:1}{Theorem~5.1},
\SGARefLink{sga3:xiii:theorem:5:1:condition:i}{(i)}\(\Rightarrow\)\SGARefLink{sga3:xiii:theorem:5:1:condition:vii}{(vii)}, applied to
\((\mathfrak h,\mathfrak d)\).

\section{Cartan subalgebras and subgroups of type \texorpdfstring{\((C)\)}{(C)}
associated with a smooth algebraic group}
\SGARefTarget{sga3:xiii:section:6}

For simplicity, in \SGARefLink{sga3:xiii:theorem:6:1}{the following theorem} we restrict ourselves to an
algebraically closed base field; the case of an arbitrary base
prescheme will be treated in the next Exposé.

\begin{theorem}
\SGARefTarget{sga3:xiii:theorem:6:1}
Let \(G\) be a smooth algebraic group over an algebraically closed
field \(k\), and let \(\mathfrak g\) be its Lie algebra. Then:
\begin{enumerate}
  \item[a)]\SGARefTarget{sga3:xiii:theorem:6:1:part:a}
  The Cartan subalgebras of \(\mathfrak g\) are conjugate.

  \item[b)]\SGARefTarget{sga3:xiii:theorem:6:1:part:b}
  Let \(\mathfrak d\) be a Cartan subalgebra of \(\mathfrak g\).
  Its normalizer \(N\) in \(G\) is smooth, and \(D=N^0\) is the
  unique smooth connected subgroup of \(G\) whose Lie algebra is
  \(\mathfrak d\). One has
  \[
    \operatorname{Norm}_G(\mathfrak d)
    =\operatorname{Norm}_G(D)
    =N,
    \qquad\text{and hence}\qquad
    D=\operatorname{Norm}_G(D)^0.
  \]

  \item[c)]\SGARefTarget{sga3:xiii:theorem:6:1:part:c}
  With \(\mathfrak d\) as in \SGARefLink{sga3:xiii:theorem:6:1:part:b}{(b)}, set, as in \SGARefLink{sga3:xiii:section:5}{Section~5},
  \[
    X=G\times^N W(\mathfrak d),
  \]
  and consider the canonical morphism
  \[
    \psi:X\longrightarrow W(\mathfrak g).
  \]
  Its fiber over \(a\in\mathfrak g\) has as its \(k\)-valued points
  the set of Cartan subalgebras of \(\mathfrak g\) containing \(a\).
  Then \(\psi\) is birational.

  \item[d)]\SGARefTarget{sga3:xiii:theorem:6:1:part:d}
  With the notation of \SGARefLink{sga3:xiii:theorem:6:1:part:c}{(c)}, let \(U\) be the largest open subset of
  \(W(\mathfrak g)\) for which \(\psi\) induces an isomorphism
  \[
    \psi^{-1}(U)\xrightarrow{\sim}U.
  \]
  Then, for \(a\in\mathfrak g\), the following conditions are
  equivalent:
  \begin{enumerate}
    \item[(i)]\SGARefTarget{sga3:xiii:theorem:6:1:part:d:condition:i}
    \(a\in U(k)\).

    \item[(ii)]\SGARefTarget{sga3:xiii:theorem:6:1:part:d:condition:ii}
    The element \(a\) belongs to one and only one Cartan subalgebra of
    \(\mathfrak g\).

    \item[(iii)]\SGARefTarget{sga3:xiii:theorem:6:1:part:d:condition:iii}
    The set of Cartan subalgebras of \(\mathfrak g\) containing \(a\)
    is finite and nonempty.

    \item[(iv)]\SGARefTarget{sga3:xiii:theorem:6:1:part:d:condition:iv}
    The fiber \(\psi^{-1}(a)\) has an isolated point.

    \item[(v)]\SGARefTarget{sga3:xiii:theorem:6:1:part:d:condition:v}
    If \(a\in\mathfrak d\), the morphism \(\psi\) is étale---or merely
    quasi-finite---at the point \((\overline e,a)\).

    \item[(vi)]\SGARefTarget{sga3:xiii:theorem:6:1:part:d:condition:vi}
    The element \(a\) is regular in \(\mathfrak g\).

    \item[(vii)]\SGARefTarget{sga3:xiii:theorem:6:1:part:d:condition:vii}
    With the notation of \SGARefLink{sga3:xiii:corollary:5:4:condition:iii}{Corollary~5.4(iii)}, one has \(N=M_a\).
  \end{enumerate}
\end{enumerate}
\end{theorem}

\begin{proof}
Apply \SGARefLink{sga3:xiii:theorem:5:1}{Theorem~5.1} with the Cartan subalgebra \(\mathfrak d\) of
\(\mathfrak g\) in place of \(\mathfrak h\). We show that the strongest
conditions, \SGARefLink{sga3:xiii:theorem:5:1:condition:xi}{(xi)} through \SGARefLink{sga3:xiii:theorem:5:1:condition:xiii}{(xiii)}, hold. This is clear either in the
form \SGARefLink{sga3:xiii:theorem:5:1:condition:xiii}{(xiii)}, by \SGARefLink{sga3:xiii:corollary:4:7}{4.7}---which says that if a regular element
\(a\in\mathfrak g\) lies in \(\mathfrak d\), then \(N=M_a\)---or in
the form
\[
  \SGARefLink{sga3:xiii:theorem:5:1:condition:xi}{(\mathrm{xi})}=\SGARefLink{sga3:xiii:theorem:5:1:condition:x}{(\mathrm{x})}+\SGARefLink{sga3:xiii:theorem:5:1:condition:i}{(\mathrm{i})}.
\]
Indeed, \SGARefLink{sga3:xiii:theorem:5:1:condition:i}{(i)} is automatic, while \SGARefLink{sga3:xiii:theorem:5:1:condition:x}{(x)} follows from the fact that a
regular point of \(\mathfrak g\) belongs to a unique Cartan subalgebra
of \(\mathfrak g\) (\SGARefLink{sga3:xiii:proposition:4:6:item:a}{4.6(a)}), and therefore, a fortiori, to a unique
conjugate of \(\mathfrak d\).

Part \SGARefLink{sga3:xiii:theorem:6:1:part:b}{(b)} now follows from \SGARefLink{sga3:xiii:corollary:5:3}{Corollary~5.3}. Part \SGARefLink{sga3:xiii:theorem:6:1:part:c}{(c)} follows because
\(\psi\) is generically étale and because a sufficiently general
point of \(\mathfrak g\)---more precisely, a regular point---belongs
to a unique Cartan subalgebra of \(\mathfrak g\). Under these
conditions, the equivalence of conditions \SGARefLink{sga3:xiii:theorem:6:1:part:d:condition:i}{(i)} through \SGARefLink{sga3:xiii:theorem:6:1:part:d:condition:v}{(v)} on \(a\)
follows immediately from Zariski's Main Theorem, applied to the
separated birational morphism \(\psi\), since \(W(\mathfrak g)\) is
normal and \(X\) is integral. The equivalence of \SGARefLink{sga3:xiii:theorem:6:1:part:d:condition:v}{(v)} and \SGARefLink{sga3:xiii:theorem:6:1:part:d:condition:vi}{(vi)} is the
special case of \SGARefLink{sga3:xiii:corollary:5:4}{Corollary~5.4},
\SGARefLink{sga3:xiii:theorem:6:1:part:d:condition:iii}{(iii)}\(\Longleftrightarrow\)\SGARefLink{sga3:xiii:theorem:6:1:part:d:condition:iv}{(iv)}, obtained after carrying \(a\) into
\(\mathfrak d\) by a suitable element \(g\in G(k)\), in view of
\SGARefLink{sga3:xiii:proposition:4:6:item:b}{4.6(b)}. Moreover, \SGARefLink{sga3:xiii:corollary:5:4}{Corollary~5.4} shows that \SGARefLink{sga3:xiii:theorem:6:1:part:d:condition:v}{(v)} and \SGARefLink{sga3:xiii:theorem:6:1:part:d:condition:vi}{(vi)} are also
equivalent to
\[
  N\supset M_a^0.
\]
By \SGARefLink{sga3:xiii:corollary:4:7}{4.7} this even implies
\(N=M_a\).\footnote{The French source prints \(N=M\) at this point; condition
\SGARefLink{sga3:xiii:theorem:6:1:part:d:condition:vii}{(vii)} and the surrounding notation require \(N=M_a\).}
This proves \SGARefLink{sga3:xiii:theorem:6:1:part:d}{(d)}.

Of course, \SGARefLink{sga3:xiii:theorem:6:1:part:b}{(b)}, \SGARefLink{sga3:xiii:theorem:6:1:part:c}{(c)}, and the equivalence of \SGARefLink{sga3:xiii:theorem:6:1:part:d:condition:i}{(i)}, \SGARefLink{sga3:xiii:theorem:6:1:part:d:condition:iv}{(iv)}, \SGARefLink{sga3:xiii:theorem:6:1:part:d:condition:v}{(v)}, \SGARefLink{sga3:xiii:theorem:6:1:part:d:condition:vi}{(vi)}, and
\SGARefLink{sga3:xiii:theorem:6:1:part:d:condition:vii}{(vii)} remain valid over an arbitrary base field. We now prove \SGARefLink{sga3:xiii:theorem:6:1:part:a}{(a)},
using the assumption that \(k\) is algebraically closed. By
\SGARefLink{sga3:xiii:theorem:5:1}{Theorem~5.1}, \SGARefLink{sga3:xiii:theorem:5:1:condition:i}{(i)}\(\Rightarrow\)\SGARefLink{sga3:xiii:theorem:5:1:condition:vii}{(vii)}, the morphism
\(\psi:X\to W(\mathfrak g)\) is dominant. Hence there is a dense open
subset \(V\) of \(W(\mathfrak g)\) such that every \(a\in V(k)\) is
the image of a point of \(X(k)\), or equivalently, belongs to a
conjugate of \(\mathfrak d\). Apply \SGARefLink{sga3:xiii:theorem:5:1:diagram:implications}{the same result} to a second Cartan
subalgebra \(\mathfrak d'\) of \(\mathfrak g\). We may choose \(V\)
so that every element of \(V(k)\) is conjugate both to an element of
\(\mathfrak d\) and to an element of \(\mathfrak d'\). Choose a
regular element of \(V(k)\). Some conjugate \(\mathfrak d''\) of
\(\mathfrak d'\) then contains an element of \(\mathfrak d\) regular
in \(\mathfrak g\), and \SGARefLink{sga3:xiii:proposition:4:6:item:a}{4.6(a)} gives
\(\mathfrak d''=\mathfrak d\). This proves \SGARefLink{sga3:xiii:theorem:6:1:part:a}{(a)} and completes the proof.
\end{proof}

\begin{definition}
\SGARefTarget{sga3:xiii:definition:6:2}
Let \(G\) be a smooth algebraic group over a field \(k\). A
\emph{subgroup of type \((C)\)} of \(G\) is a smooth connected
algebraic subgroup whose Lie algebra is a Cartan subalgebra of
\(\mathfrak g\). The \emph{infinitesimal rank} of \(G\) is the
nilpotent rank of its Lie algebra \(\mathfrak g\); it is equal to the
dimension of every subgroup of type \((C)\) of \(G\).
\end{definition}

\SGARefLink{sga3:xiii:theorem:6:1:part:b}{Theorem~6.1(b)} immediately gives:

\begin{corollary}
\SGARefTarget{sga3:xiii:corollary:6:3}
Under the hypotheses of \SGARefLink{sga3:xiii:definition:6:2}{Definition~6.2}, the map
\[
  D\longmapsto\mathfrak d=\operatorname{Lie}(D)
\]
establishes a bijection between the subgroups of type \((C)\) of
\(G\) and the Cartan subalgebras \(\mathfrak d\) of \(\mathfrak g\).
If \(D\) is a subgroup of type \((C)\) of \(G\), then \(D\) is its
own connected normalizer:
\[
  D=\operatorname{Norm}_G(D)^0.
\]
\end{corollary}

Combining \SGARefLink{sga3:xiii:theorem:6:1:part:a}{Theorem~6.1(a)} and \SGARefLink{sga3:xiii:corollary:6:3}{Corollary~6.3} gives:

\begin{corollary}
\SGARefTarget{sga3:xiii:corollary:6:4}
If \(k\) is algebraically closed, the subgroups of type \((C)\) of
\(G\) are conjugate to one another.
\end{corollary}

\begin{corollary}
\SGARefTarget{sga3:xiii:corollary:6:5}
Let \(G\) be a smooth algebraic group over an algebraically closed
field \(k\), and let \(H\) be a smooth algebraic subgroup of \(G\),
with Lie algebras \(\mathfrak g\) and \(\mathfrak h\), respectively.
Then \(\mathfrak h\) contains a Cartan subalgebra of \(\mathfrak g\)
if and only if \(H\) contains a subgroup \(D\) of type \((C)\) of
\(G\).
\end{corollary}

The condition is plainly sufficient. It is also necessary, since by
\SGARefLink{sga3:xiii:corollary:5:5}{Corollary~5.5},
\[
  H\supset D
  \quad\Longleftrightarrow\quad
  \mathfrak h\supset\mathfrak d.
\]

\refstepcounter{theorem}
\par\medskip
\noindent\SGARefTarget{sga3:xiii:remarks:6:6}%
\textbf{Remarks \thetheorem.}
\begin{enumerate}
  \item[a)]\SGARefTarget{sga3:xiii:remarks:6:6:part:a}
  The connected subgroups \(H\) of \(G\) described in
  \SGARefLink{sga3:xiii:corollary:6:5}{Corollary~6.5} correspond bijectively to the Lie subalgebras
  \(\mathfrak h\) of \(\mathfrak g\) satisfying the strongest
  condition, \SGARefLink{sga3:xiii:theorem:5:1:condition:xii}{(xii)}, of \SGARefLink{sga3:xiii:theorem:5:1}{Theorem~5.1}. By \SGARefLink{sga3:xiii:corollary:5:5}{Corollary~5.5}, inclusion
  relations among such subgroups can already be read from their Lie
  algebras.

  \item[b)]\SGARefTarget{sga3:xiii:remarks:6:6:part:b}
  Continue to suppose that \(k\) is algebraically closed, and let
  \(D\) be a subgroup of type \((C)\) of \(G\). It is then easy to
  show that \(D\) contains a Cartan subgroup \(C\) of \(G\). Indeed,
  put \(V=\mathfrak g/\mathfrak d\). For a general element
  \(a\in\mathfrak d\), the endomorphism
  \(\operatorname{ad}(a)\) of \(V\) is injective; it suffices to take
  \(a\) regular in \(\mathfrak g\). From this one readily deduces that,
  for a general \(g\in D(k)\), the action of
  \(\operatorname{Ad}(g)\) on \(V\) has no nonzero invariant vector.
  Thus \SGARefLink{sga3:xiii:theorem:2:1}{2.1}, \SGARefLink{sga3:xiii:theorem:2:1:item:vii}{(vii)}\(\Rightarrow\)\SGARefLink{sga3:xiii:theorem:2:1:item:i}{(i)}, applies. In fact, in the next
  Exposé we shall prove the more precise result that every Cartan
  subgroup \(C\) of \(G\) is contained in one and only one subgroup
  of type \((C)\) of \(G\).

  \item[c)]\SGARefTarget{sga3:xiii:remarks:6:6:part:c}
  In general one must not confuse Cartan subgroups with subgroups of
  type \((C)\). The subgroups of type \((C)\) of \(G\) are Cartan
  subgroups if and only if they are nilpotent, since Cartan subgroups
  are maximal nilpotent subgroups. It may happen that \(\mathfrak g\)
  is nilpotent although \(G\) is not; an example is
  \(\operatorname{SL}(2)_k\) when \(k\) has characteristic \(2\).
  Then the Cartan subalgebras of \(\mathfrak g\) are all equal to
  \(\mathfrak g\). Thus, if \(G\) is connected, \(G\) is a subgroup
  of type \((C)\) of itself, but it is not a Cartan subgroup of
  itself.

  On the other hand, we shall see in Exposé~XIV that if \(G\) is an
  adjoint semisimple group, its subgroups of type \((C)\) are precisely
  its Cartan subgroups. The same conclusion holds in characteristic
  zero for every smooth algebraic group \(G\) over \(k\). This follows
  at once from the fact that, in characteristic zero, a connected
  algebraic group is nilpotent if and only if its Lie algebra is
  nilpotent. In turn, this follows from the fact that in
  characteristic zero the center \(Z\) of a connected algebraic group
  \(G\) has Lie algebra equal to the center of \(\mathfrak g\).
  Indeed, one obtains a priori the space of invariants in
  \(\mathfrak g\) under the adjoint representation of \(G\). Since
  \(G\) is connected and \(k\) has characteristic zero, these are also
  the infinitesimal invariants of the adjoint representation of
  \(\mathfrak g\), namely the center of \(\mathfrak g\). Finally,
  Cartier's theorem (cf.\ VI\textsubscript{B}~1.6.1) shows that \(Z\)
  is smooth.
\end{enumerate}

% END inputs/XIII/source/tex/components/04_expose_XIII_sections5_6.tex


